
\documentclass{article}
\usepackage{colortbl}
\usepackage{makecell}
\usepackage{multirow}
\usepackage{supertabular}

\begin{document}

\newcounter{utterance}

\centering \large Interaction Transcript for game `clean\_up', experiment `0\_easy\_7obj\_it', episode 1 with deepseek{-}v4{-}pro.
\vspace{24pt}

{ \footnotesize  \setcounter{utterance}{1}
\setlength{\tabcolsep}{0pt}
\begin{supertabular}{c@{$\;$}|p{.15\linewidth}@{}p{.15\linewidth}p{.15\linewidth}p{.15\linewidth}p{.15\linewidth}p{.15\linewidth}}
    \# & $\;$A & \multicolumn{4}{c}{Game Master} & $\;\:$B\\
    \hline

    \theutterance \stepcounter{utterance}  
    & & \multicolumn{4}{p{0.6\linewidth}}{
        \cellcolor[rgb]{0.9,0.9,0.9}{
            \makecell[{{p{\linewidth}}}]{
                \texttt{\tiny{[P1$\langle$GM]}}
                \texttt{Sono il tuo maestro di gioco e stai giocando un gioco collaborativo con la seguente griglia come tabellone di gioco:} \\
\\ 
\texttt{```} \\
\texttt{ 1234567} \\
\texttt{╔═══════╗} \\
\texttt{║◌◌◌◌◌T◌║ 1} \\
\texttt{║P◌◌◌◌◌◌║ 2} \\
\texttt{║◌◌AD◌┌─╢ 3} \\
\texttt{║◌◌◌┌─┤S║ 4} \\
\texttt{║O◌◌│◌├─╢ 5} \\
\texttt{╟───┼─┘◌║ 6} \\
\texttt{║◌M◌│◌◌◌║ 7} \\
\texttt{╚═══╧═══╝} \\
\\ 
\texttt{```} \\
\\ 
\texttt{* Il bordo superiore mostra le coordinate x che aumentano verso destra, e il bordo destro le coordinate y che aumentano verso il basso.} \\
\texttt{* I seguenti oggetti sono posizionati casualmente sulla tua griglia: 'P', 'O', 'T', 'S', 'D', 'A', 'M'.} \\
\\ 
\\ 
\texttt{L'altro giocatore vede una variante del tabellone di gioco, dove gli oggetti sono posizionati in diverse posizioni casuali. Non puoi vedere il tabellone dell'altro giocatore e lui non può vedere il tuo.} \\
\\ 
\\ 
\texttt{**Obiettivo:**} \\
\\ 
\texttt{Entrambi i giocatori devono spostare gli oggetti sui rispettivi sfondi in modo che oggetti identici finiscano alle stesse coordinate. Devi comunicare con l'altro giocatore per concordare una configurazione obiettivo comune.} \\
\\ 
\\ 
\texttt{**Regole:**} \\
\\ 
\texttt{* In ogni turno, puoi inviare esattamente uno dei seguenti due comandi:} \\
\texttt{ 1. `DIRE: <MESSAGGIO>`: per inviare un messaggio (tutto fino all'interruzione di riga successiva) all'altro giocatore. Lo inoltrerò al tuo partner.} \\
\texttt{2. `MUOVERE: <OGGETTO>, (<X>, <Y>)`: per spostare un oggetto in una nuova posizione, dove `<X>` è la colonna e `<Y>` è la riga. Ti informerò se il tuo spostamento ha avuto successo o meno.} \\
\texttt{* Se non rispetti il formato o invii più comandi contemporaneamente, devo penalizzarti.} \\
\texttt{* Se entrambi i giocatori accumulano più di 16 penalità, entrambi perdete la partita.} \\
\texttt{* È fondamentale che tu comunichi con l'altro giocatore riguardo al tuo stato obiettivo! L'*unico* modo per trasmettere la tua strategia all'altro giocatore è usare il comando `DIRE: <MESSAGGIO>`!} \\
\\ 
\\ 
\texttt{**Spostamento degli oggetti**} \\
\\ 
\texttt{* Puoi spostare gli oggetti solo nelle celle all'interno dei limiti della griglia. La cella di destinazione deve essere vuota, cioè deve contenere solo il simbolo '◌'.} \\
\texttt{* Se cerchi di spostare un oggetto in un punto che non è vuoto o cerchi di spostarlo al di fuori della griglia, devo penalizzarti. Avrai un altro tentativo.} \\
\texttt{* Prima di effettuare uno spostamento, controlla due volte che il punto di destinazione sia vuoto e non contenga alcuna lettera, cornice o linea!} \\
\\ 
\\ 
\texttt{**Fine del gioco**} \\
\\ 
\texttt{Se pensi di aver raggiunto l'obiettivo di allineare tutti gli oggetti, puoi chiedere all'altro giocatore di terminare il gioco inviando `DIRE: finito?`. Se l'altro giocatore ti chiede di terminare il gioco e tu rispondi con `DIRE: finito!`, il gioco terminerà.} \\
\\ 
\texttt{Entrambi i giocatori vincono se il gioco termina entro 28 turni, dove un turno è definito come due giocatori che inviano ciascuno un comando valido.} \\
\\ 
\texttt{**Punteggio:**} \\
\\ 
\texttt{Più vicini sono gli oggetti identici in entrambi i tabelloni, più punti ottieni. Le penalità riducono i tuoi punti. Riuscirai a battere il record?} \\
\\ 
\\ 
\texttt{Invia un messaggio all'altro giocatore per iniziare la partita!} \\
            }
        }
    }
    & \\ \\

    \theutterance \stepcounter{utterance}  
    & \multicolumn{4}{p{0.6\linewidth}}{
        \cellcolor[rgb]{0.9,0.9,0.9}{
            \makecell[{{p{\linewidth}}}]{
                \texttt{\tiny{[P1$\rangle$GM]}}
                \texttt{DIRE: Ciao! Iniziamo a confrontare le posizioni. Ecco la mia configurazione:} \\
\texttt{P: (1,2)} \\
\texttt{O: (1,5)} \\
\texttt{T: (6,1)} \\
\texttt{S: (6,4)} \\
\texttt{D: (3,3)} \\
\texttt{A: (2,3)} \\
\texttt{M: (2,7)} \\
\texttt{Quali sono le tue posizioni?} \\
            }
        }
    }
    & & \\ \\

    \theutterance \stepcounter{utterance}  
    & & & \multicolumn{2}{p{0.3\linewidth}}{
        \cellcolor[rgb]{0.9,0.9,0.9}{
            \makecell[{{p{\linewidth}}}]{
                \texttt{\tiny{[GM$|$GM]}}
                \texttt{DIRE: Ciao! Iniziamo a confrontare le posizioni. Ecco la mia configurazione:} \\
\texttt{P: (1,2)} \\
\texttt{O: (1,5)} \\
\texttt{T: (6,1)} \\
\texttt{S: (6,4)} \\
\texttt{D: (3,3)} \\
\texttt{A: (2,3)} \\
\texttt{M: (2,7)} \\
\texttt{Quali sono le tue posizioni?} \\
            }
        }
    }
    & & \\ \\

    \theutterance \stepcounter{utterance}  
    & & & \multicolumn{2}{p{0.3\linewidth}}{
        \cellcolor[rgb]{0.9,0.9,0.9}{
            \makecell[{{p{\linewidth}}}]{
                \texttt{\tiny{[GM$|$GM]}}
                \texttt{Invalid format: tail is not empty: '} \\
\texttt{P: (1,2)} \\
\texttt{O: (1,5)} \\
\texttt{T: (6,1)} \\
\texttt{S: (6,4)} \\
\texttt{D: (3,3)} \\
\texttt{A: (2,3)} \\
\texttt{M: (2,7)} \\
\texttt{Quali sono le tue posizioni?} \\
            }
        }
    }
    & & \\ \\

    \theutterance \stepcounter{utterance}  
    & & \multicolumn{4}{p{0.6\linewidth}}{
        \cellcolor[rgb]{0.9,0.9,0.9}{
            \makecell[{{p{\linewidth}}}]{
                \texttt{\tiny{[P1$\langle$GM]}}
                \texttt{Penalità: Il tuo messaggio non deve contenere nulla dopo il comando!} \\
\\ 
\texttt{Assicurati che la tua risposta contenga solo `DIRE: <MESSAGGIO>` o `MUOVERE: <OGGETTO>, (<X>, <Y>)`, e nient'altro!} \\
\\ 
\texttt{Avete accumulato collettivamente 1 di 16 penalità.} \\
\texttt{Per favore riprova!} \\
            }
        }
    }
    & \\ \\

    \theutterance \stepcounter{utterance}  
    & \multicolumn{4}{p{0.6\linewidth}}{
        \cellcolor[rgb]{0.9,0.9,0.9}{
            \makecell[{{p{\linewidth}}}]{
                \texttt{\tiny{[P1$\rangle$GM]}}
                \texttt{DIRE: Ciao! Iniziamo a confrontare le posizioni. Ecco la mia configurazione: P a (1,2), O a (1,5), T a (6,1), S a (6,4), D a (3,3), A a (2,3), M a (2,7). Quali sono le tue posizioni?} \\
            }
        }
    }
    & & \\ \\

    \theutterance \stepcounter{utterance}  
    & & & \multicolumn{2}{p{0.3\linewidth}}{
        \cellcolor[rgb]{0.9,0.9,0.9}{
            \makecell[{{p{\linewidth}}}]{
                \texttt{\tiny{[GM$|$GM]}}
                \texttt{DIRE: Ciao! Iniziamo a confrontare le posizioni. Ecco la mia configurazione: P a (1,2), O a (1,5), T a (6,1), S a (6,4), D a (3,3), A a (2,3), M a (2,7). Quali sono le tue posizioni?} \\
            }
        }
    }
    & & \\ \\

    \theutterance \stepcounter{utterance}  
    & & \multicolumn{4}{p{0.6\linewidth}}{
        \cellcolor[rgb]{0.9,0.9,0.9}{
            \makecell[{{p{\linewidth}}}]{
                \texttt{\tiny{[GM$\rangle$P2]}}
                \texttt{Sono il tuo maestro di gioco e stai giocando un gioco collaborativo con la seguente griglia come tabellone di gioco:} \\
\\ 
\texttt{```} \\
\texttt{ 1234567} \\
\texttt{╔═══════╗} \\
\texttt{║◌◌◌◌◌S◌║ 1} \\
\texttt{║◌◌◌◌O◌◌║ 2} \\
\texttt{║◌◌◌M◌┌─╢ 3} \\
\texttt{║◌◌A┌─┤◌║ 4} \\
\texttt{║◌◌◌│◌├─╢ 5} \\
\texttt{╟───┼─┘◌║ 6} \\
\texttt{║D◌T│P◌◌║ 7} \\
\texttt{╚═══╧═══╝} \\
\\ 
\texttt{```} \\
\\ 
\texttt{* Il bordo superiore mostra le coordinate x che aumentano verso destra, e il bordo destro le coordinate y che aumentano verso il basso.} \\
\texttt{* I seguenti oggetti sono posizionati casualmente sulla tua griglia: 'P', 'O', 'T', 'S', 'D', 'A', 'M'.} \\
\\ 
\\ 
\texttt{L'altro giocatore vede una variante del tabellone di gioco, dove gli oggetti sono posizionati in diverse posizioni casuali. Non puoi vedere il tabellone dell'altro giocatore e lui non può vedere il tuo.} \\
\\ 
\\ 
\texttt{**Obiettivo:**} \\
\\ 
\texttt{Entrambi i giocatori devono spostare gli oggetti sui rispettivi sfondi in modo che oggetti identici finiscano alle stesse coordinate. Devi comunicare con l'altro giocatore per concordare una configurazione obiettivo comune.} \\
\\ 
\\ 
\texttt{**Regole:**} \\
\\ 
\texttt{* In ogni turno, puoi inviare esattamente uno dei seguenti due comandi:} \\
\texttt{ 1. `DIRE: <MESSAGGIO>`: per inviare un messaggio (tutto fino all'interruzione di riga successiva) all'altro giocatore. Lo inoltrerò al tuo partner.} \\
\texttt{2. `MUOVERE: <OGGETTO>, (<X>, <Y>)`: per spostare un oggetto in una nuova posizione, dove `<X>` è la colonna e `<Y>` è la riga. Ti informerò se il tuo spostamento ha avuto successo o meno.} \\
\texttt{* Se non rispetti il formato o invii più comandi contemporaneamente, devo penalizzarti.} \\
\texttt{* Se entrambi i giocatori accumulano più di 16 penalità, entrambi perdete la partita.} \\
\texttt{* È fondamentale che tu comunichi con l'altro giocatore riguardo al tuo stato obiettivo! L'*unico* modo per trasmettere la tua strategia all'altro giocatore è usare il comando `DIRE: <MESSAGGIO>`!} \\
\\ 
\\ 
\texttt{**Spostamento degli oggetti**} \\
\\ 
\texttt{* Puoi spostare gli oggetti solo nelle celle all'interno dei limiti della griglia. La cella di destinazione deve essere vuota, cioè deve contenere solo il simbolo '◌'.} \\
\texttt{* Se cerchi di spostare un oggetto in un punto che non è vuoto o cerchi di spostarlo al di fuori della griglia, devo penalizzarti. Avrai un altro tentativo.} \\
\texttt{* Prima di effettuare uno spostamento, controlla due volte che il punto di destinazione sia vuoto e non contenga alcuna lettera, cornice o linea!} \\
\\ 
\\ 
\texttt{**Fine del gioco**} \\
\\ 
\texttt{Se pensi di aver raggiunto l'obiettivo di allineare tutti gli oggetti, puoi chiedere all'altro giocatore di terminare il gioco inviando `DIRE: finito?`. Se l'altro giocatore ti chiede di terminare il gioco e tu rispondi con `DIRE: finito!`, il gioco terminerà.} \\
\\ 
\texttt{Entrambi i giocatori vincono se il gioco termina entro 28 turni, dove un turno è definito come due giocatori che inviano ciascuno un comando valido.} \\
\\ 
\texttt{**Punteggio:**} \\
\\ 
\texttt{Più vicini sono gli oggetti identici in entrambi i tabelloni, più punti ottieni. Le penalità riducono i tuoi punti. Riuscirai a battere il record?} \\
\\ 
\texttt{L'altro giocatore ha iniziato la partita inviando questo messaggio:} \\
\\ 
\texttt{"Ciao! Iniziamo a confrontare le posizioni. Ecco la mia configurazione: P a (1,2), O a (1,5), T a (6,1), S a (6,4), D a (3,3), A a (2,3), M a (2,7). Quali sono le tue posizioni?"} \\
\\ 
\texttt{Qual è il tuo primo comando?} \\
            }
        }
    }
    & \\ \\

    \theutterance \stepcounter{utterance}  
    & & & \multicolumn{4}{p{0.6\linewidth}}{
        \cellcolor[rgb]{0.9,0.9,0.9}{
            \makecell[{{p{\linewidth}}}]{
                \texttt{\tiny{[GM$\langle$P2]}}
                \texttt{Analizzo la mia griglia per determinare le posizioni dei miei oggetti:} \\
\\ 
\texttt{**Mia configurazione:**} \\
\texttt{{-} **P**: (5,7)} \\
\texttt{{-} **O**: (5,2)} \\
\texttt{{-} **T**: (6,7)} \\
\texttt{{-} **S**: (6,1)} \\
\texttt{{-} **D**: (1,7)} \\
\texttt{{-} **A**: (3,4)} \\
\texttt{{-} **M**: (4,3)} \\
\\ 
\texttt{**Configurazione dell'altro giocatore:**} \\
\texttt{{-} P: (1,2)} \\
\texttt{{-} O: (1,5)} \\
\texttt{{-} T: (6,1)} \\
\texttt{{-} S: (6,4)} \\
\texttt{{-} D: (3,3)} \\
\texttt{{-} A: (2,3)} \\
\texttt{{-} M: (2,7)} \\
\\ 
\texttt{Ora confronto le posizioni. L'obiettivo è trovare una configurazione comune. Mi conviene comunicare le mie posizioni per confrontarle.} \\
\\ 
\texttt{**Primo comando:**} \\
\texttt{```} \\
\texttt{DIRE: Ecco le mie posizioni: P a (5,7), O a (5,2), T a (6,7), S a (6,1), D a (1,7), A a (3,4), M a (4,3). Confrontiamo e decidiamo dove spostare ogni oggetto.} \\
\texttt{```} \\
            }
        }
    }
     \\ \\

    \theutterance \stepcounter{utterance}  
    & & & \multicolumn{2}{p{0.3\linewidth}}{
        \cellcolor[rgb]{0.9,0.9,0.9}{
            \makecell[{{p{\linewidth}}}]{
                \texttt{\tiny{[GM$|$GM]}}
                \texttt{Analizzo la mia griglia per determinare le posizioni dei miei oggetti:} \\
\\ 
\texttt{**Mia configurazione:**} \\
\texttt{{-} **P**: (5,7)} \\
\texttt{{-} **O**: (5,2)} \\
\texttt{{-} **T**: (6,7)} \\
\texttt{{-} **S**: (6,1)} \\
\texttt{{-} **D**: (1,7)} \\
\texttt{{-} **A**: (3,4)} \\
\texttt{{-} **M**: (4,3)} \\
\\ 
\texttt{**Configurazione dell'altro giocatore:**} \\
\texttt{{-} P: (1,2)} \\
\texttt{{-} O: (1,5)} \\
\texttt{{-} T: (6,1)} \\
\texttt{{-} S: (6,4)} \\
\texttt{{-} D: (3,3)} \\
\texttt{{-} A: (2,3)} \\
\texttt{{-} M: (2,7)} \\
\\ 
\texttt{Ora confronto le posizioni. L'obiettivo è trovare una configurazione comune. Mi conviene comunicare le mie posizioni per confrontarle.} \\
\\ 
\texttt{**Primo comando:**} \\
\texttt{```} \\
\texttt{DIRE: Ecco le mie posizioni: P a (5,7), O a (5,2), T a (6,7), S a (6,1), D a (1,7), A a (3,4), M a (4,3). Confrontiamo e decidiamo dove spostare ogni oggetto.} \\
\texttt{```} \\
            }
        }
    }
    & & \\ \\

    \theutterance \stepcounter{utterance}  
    & & & \multicolumn{2}{p{0.3\linewidth}}{
        \cellcolor[rgb]{0.9,0.9,0.9}{
            \makecell[{{p{\linewidth}}}]{
                \texttt{\tiny{[GM$|$GM]}}
                \texttt{Invalid format: head is not empty: 'Analizzo la mia griglia per determinare le posizioni dei miei oggetti:} \\
\\ 
\texttt{**Mia configurazione:**} \\
\texttt{{-} **P**: (5,7)} \\
\texttt{{-} **O**: (5,2)} \\
\texttt{{-} **T**: (6,7)} \\
\texttt{{-} **S**: (6,1)} \\
\texttt{{-} **D**: (1,7)} \\
\texttt{{-} **A**: (3,4)} \\
\texttt{{-} **M**: (4,3)} \\
\\ 
\texttt{**Configurazione dell'altro giocatore:**} \\
\texttt{{-} P: (1,2)} \\
\texttt{{-} O: (1,5)} \\
\texttt{{-} T: (6,1)} \\
\texttt{{-} S: (6,4)} \\
\texttt{{-} D: (3,3)} \\
\texttt{{-} A: (2,3)} \\
\texttt{{-} M: (2,7)} \\
\\ 
\texttt{Ora confronto le posizioni. L'obiettivo è trovare una configurazione comune. Mi conviene comunicare le mie posizioni per confrontarle.} \\
\\ 
\texttt{**Primo comando:**} \\
\\ 
\texttt{'} \\
            }
        }
    }
    & & \\ \\

    \theutterance \stepcounter{utterance}  
    & & \multicolumn{4}{p{0.6\linewidth}}{
        \cellcolor[rgb]{0.9,0.9,0.9}{
            \makecell[{{p{\linewidth}}}]{
                \texttt{\tiny{[GM$\rangle$P2]}}
                \texttt{Penalità: Il tuo messaggio non deve contenere nulla prima del comando!} \\
\\ 
\texttt{Assicurati che la tua risposta contenga solo `DIRE: <MESSAGGIO>` o `MUOVERE: <OGGETTO>, (<X>, <Y>)`, e nient'altro!} \\
\\ 
\texttt{Avete accumulato collettivamente 2 di 16 penalità.} \\
\texttt{Per favore riprova!} \\
            }
        }
    }
    & \\ \\

    \theutterance \stepcounter{utterance}  
    & & & \multicolumn{4}{p{0.6\linewidth}}{
        \cellcolor[rgb]{0.9,0.9,0.9}{
            \makecell[{{p{\linewidth}}}]{
                \texttt{\tiny{[GM$\langle$P2]}}
                \texttt{DIRE: Ecco le mie posizioni: P a (5,7), O a (5,2), T a (6,7), S a (6,1), D a (1,7), A a (3,4), M a (4,3). Confrontiamo e decidiamo dove spostare ogni oggetto.} \\
            }
        }
    }
     \\ \\

    \theutterance \stepcounter{utterance}  
    & & & \multicolumn{2}{p{0.3\linewidth}}{
        \cellcolor[rgb]{0.9,0.9,0.9}{
            \makecell[{{p{\linewidth}}}]{
                \texttt{\tiny{[GM$|$GM]}}
                \texttt{DIRE: Ecco le mie posizioni: P a (5,7), O a (5,2), T a (6,7), S a (6,1), D a (1,7), A a (3,4), M a (4,3). Confrontiamo e decidiamo dove spostare ogni oggetto.} \\
            }
        }
    }
    & & \\ \\

    \theutterance \stepcounter{utterance}  
    & & \multicolumn{4}{p{0.6\linewidth}}{
        \cellcolor[rgb]{0.9,0.9,0.9}{
            \makecell[{{p{\linewidth}}}]{
                \texttt{\tiny{[P1$\langle$GM]}}
                \texttt{Il tuo messaggio è stato trasmesso all'altro giocatore.} \\
\\ 
\\ 
\texttt{Stai giocando il round 1 di un massimo di 28.} \\
\texttt{Avete accumulato collettivamente 2 di 16 penalità.} \\
\texttt{L'altro giocatore ha inviato questo messaggio:} \\
\\ 
\\ 
\texttt{"Ecco le mie posizioni: P a (5,7), O a (5,2), T a (6,7), S a (6,1), D a (1,7), A a (3,4), M a (4,3). Confrontiamo e decidiamo dove spostare ogni oggetto."} \\
\\ 
\\ 
\texttt{Qual è il tuo prossimo comando?} \\
            }
        }
    }
    & \\ \\

    \theutterance \stepcounter{utterance}  
    & \multicolumn{4}{p{0.6\linewidth}}{
        \cellcolor[rgb]{0.9,0.9,0.9}{
            \makecell[{{p{\linewidth}}}]{
                \texttt{\tiny{[P1$\rangle$GM]}}
                \texttt{DIRE: Ecco un possibile obiettivo comune. Propongo di spostare: P a (3,4) ma è occupata da A? Aspetta, controlliamo una cella vuota. Propongo: P a (4,4), O a (3,2), T a (6,4), S a (6,3), D a (2,4), A a (3,5), M a (3,6). Che ne pensi?} \\
            }
        }
    }
    & & \\ \\

    \theutterance \stepcounter{utterance}  
    & & & \multicolumn{2}{p{0.3\linewidth}}{
        \cellcolor[rgb]{0.9,0.9,0.9}{
            \makecell[{{p{\linewidth}}}]{
                \texttt{\tiny{[GM$|$GM]}}
                \texttt{DIRE: Ecco un possibile obiettivo comune. Propongo di spostare: P a (3,4) ma è occupata da A? Aspetta, controlliamo una cella vuota. Propongo: P a (4,4), O a (3,2), T a (6,4), S a (6,3), D a (2,4), A a (3,5), M a (3,6). Che ne pensi?} \\
            }
        }
    }
    & & \\ \\

    \theutterance \stepcounter{utterance}  
    & & \multicolumn{4}{p{0.6\linewidth}}{
        \cellcolor[rgb]{0.9,0.9,0.9}{
            \makecell[{{p{\linewidth}}}]{
                \texttt{\tiny{[GM$\rangle$P2]}}
                \texttt{Il tuo messaggio è stato trasmesso all'altro giocatore.} \\
\\ 
\\ 
\texttt{Stai giocando il round 2 di un massimo di 28.} \\
\texttt{Avete accumulato collettivamente 2 di 16 penalità.} \\
\texttt{L'altro giocatore ha inviato questo messaggio:} \\
\\ 
\\ 
\texttt{"Ecco un possibile obiettivo comune. Propongo di spostare: P a (3,4) ma è occupata da A? Aspetta, controlliamo una cella vuota. Propongo: P a (4,4), O a (3,2), T a (6,4), S a (6,3), D a (2,4), A a (3,5), M a (3,6). Che ne pensi?"} \\
\\ 
\\ 
\texttt{Qual è il tuo prossimo comando?} \\
            }
        }
    }
    & \\ \\

    \theutterance \stepcounter{utterance}  
    & & & \multicolumn{4}{p{0.6\linewidth}}{
        \cellcolor[rgb]{0.9,0.9,0.9}{
            \makecell[{{p{\linewidth}}}]{
                \texttt{\tiny{[GM$\langle$P2]}}
                \texttt{Controllo le posizioni proposte sulla mia griglia:} \\
\\ 
\texttt{{-} **P a (4,4)**: attualmente c'è una linea `┌` – non è vuoto.} \\
\texttt{{-} **O a (3,2)**: vuoto `◌` – OK.} \\
\texttt{{-} **T a (6,4)**: vuoto `◌` – OK.} \\
\texttt{{-} **S a (6,3)**: c'è `─` – non è vuoto.} \\
\texttt{{-} **D a (2,4)**: vuoto `◌` – OK.} \\
\texttt{{-} **A a (3,5)**: vuoto `◌` – OK.} \\
\texttt{{-} **M a (3,6)**: c'è `─` – non è vuoto.} \\
\\ 
\texttt{Ci sono problemi con alcune celle. Devo comunicarlo e proporre alternative.} \\
\\ 
\texttt{```} \\
\texttt{DIRE: P a (4,4) non è vuoto (c'è una linea), S a (6,3) e M a (3,6) anche. Propongo alternative: P a (4,4) spostiamolo a (4,5), S a (6,3) a (7,3), M a (3,6) a (3,7). Per il resto OK: O a (3,2), T a (6,4), D a (2,4), A a (3,5).} \\
\texttt{```} \\
            }
        }
    }
     \\ \\

    \theutterance \stepcounter{utterance}  
    & & & \multicolumn{2}{p{0.3\linewidth}}{
        \cellcolor[rgb]{0.9,0.9,0.9}{
            \makecell[{{p{\linewidth}}}]{
                \texttt{\tiny{[GM$|$GM]}}
                \texttt{Controllo le posizioni proposte sulla mia griglia:} \\
\\ 
\texttt{{-} **P a (4,4)**: attualmente c'è una linea `┌` – non è vuoto.} \\
\texttt{{-} **O a (3,2)**: vuoto `◌` – OK.} \\
\texttt{{-} **T a (6,4)**: vuoto `◌` – OK.} \\
\texttt{{-} **S a (6,3)**: c'è `─` – non è vuoto.} \\
\texttt{{-} **D a (2,4)**: vuoto `◌` – OK.} \\
\texttt{{-} **A a (3,5)**: vuoto `◌` – OK.} \\
\texttt{{-} **M a (3,6)**: c'è `─` – non è vuoto.} \\
\\ 
\texttt{Ci sono problemi con alcune celle. Devo comunicarlo e proporre alternative.} \\
\\ 
\texttt{```} \\
\texttt{DIRE: P a (4,4) non è vuoto (c'è una linea), S a (6,3) e M a (3,6) anche. Propongo alternative: P a (4,4) spostiamolo a (4,5), S a (6,3) a (7,3), M a (3,6) a (3,7). Per il resto OK: O a (3,2), T a (6,4), D a (2,4), A a (3,5).} \\
\texttt{```} \\
            }
        }
    }
    & & \\ \\

    \theutterance \stepcounter{utterance}  
    & & & \multicolumn{2}{p{0.3\linewidth}}{
        \cellcolor[rgb]{0.9,0.9,0.9}{
            \makecell[{{p{\linewidth}}}]{
                \texttt{\tiny{[GM$|$GM]}}
                \texttt{Invalid format: head is not empty: 'Controllo le posizioni proposte sulla mia griglia:} \\
\\ 
\texttt{{-} **P a (4,4)**: attualmente c'è una linea ┌ – non è vuoto.} \\
\texttt{{-} **O a (3,2)**: vuoto ◌ – OK.} \\
\texttt{{-} **T a (6,4)**: vuoto ◌ – OK.} \\
\texttt{{-} **S a (6,3)**: c'è ─ – non è vuoto.} \\
\texttt{{-} **D a (2,4)**: vuoto ◌ – OK.} \\
\texttt{{-} **A a (3,5)**: vuoto ◌ – OK.} \\
\texttt{{-} **M a (3,6)**: c'è ─ – non è vuoto.} \\
\\ 
\texttt{Ci sono problemi con alcune celle. Devo comunicarlo e proporre alternative.} \\
\\ 
\\ 
\texttt{'} \\
            }
        }
    }
    & & \\ \\

    \theutterance \stepcounter{utterance}  
    & & \multicolumn{4}{p{0.6\linewidth}}{
        \cellcolor[rgb]{0.9,0.9,0.9}{
            \makecell[{{p{\linewidth}}}]{
                \texttt{\tiny{[GM$\rangle$P2]}}
                \texttt{Penalità: Il tuo messaggio non deve contenere nulla prima del comando!} \\
\\ 
\texttt{Assicurati che la tua risposta contenga solo `DIRE: <MESSAGGIO>` o `MUOVERE: <OGGETTO>, (<X>, <Y>)`, e nient'altro!} \\
\\ 
\texttt{Avete accumulato collettivamente 3 di 16 penalità.} \\
\texttt{Per favore riprova!} \\
            }
        }
    }
    & \\ \\

    \theutterance \stepcounter{utterance}  
    & & & \multicolumn{4}{p{0.6\linewidth}}{
        \cellcolor[rgb]{0.9,0.9,0.9}{
            \makecell[{{p{\linewidth}}}]{
                \texttt{\tiny{[GM$\langle$P2]}}
                \texttt{DIRE: P a (4,4) non è vuoto (c'è una linea), S a (6,3) e M a (3,6) anche. Propongo alternative: P a (4,5), S a (7,3), M a (3,7). Per il resto OK: O a (3,2), T a (6,4), D a (2,4), A a (3,5).} \\
            }
        }
    }
     \\ \\

    \theutterance \stepcounter{utterance}  
    & & & \multicolumn{2}{p{0.3\linewidth}}{
        \cellcolor[rgb]{0.9,0.9,0.9}{
            \makecell[{{p{\linewidth}}}]{
                \texttt{\tiny{[GM$|$GM]}}
                \texttt{DIRE: P a (4,4) non è vuoto (c'è una linea), S a (6,3) e M a (3,6) anche. Propongo alternative: P a (4,5), S a (7,3), M a (3,7). Per il resto OK: O a (3,2), T a (6,4), D a (2,4), A a (3,5).} \\
            }
        }
    }
    & & \\ \\

    \theutterance \stepcounter{utterance}  
    & & \multicolumn{4}{p{0.6\linewidth}}{
        \cellcolor[rgb]{0.9,0.9,0.9}{
            \makecell[{{p{\linewidth}}}]{
                \texttt{\tiny{[P1$\langle$GM]}}
                \texttt{Il tuo messaggio è stato trasmesso all'altro giocatore.} \\
\\ 
\\ 
\texttt{Stai giocando il round 2 di un massimo di 28.} \\
\texttt{Avete accumulato collettivamente 3 di 16 penalità.} \\
\texttt{L'altro giocatore ha inviato questo messaggio:} \\
\\ 
\\ 
\texttt{"P a (4,4) non è vuoto (c'è una linea), S a (6,3) e M a (3,6) anche. Propongo alternative: P a (4,5), S a (7,3), M a (3,7). Per il resto OK: O a (3,2), T a (6,4), D a (2,4), A a (3,5)."} \\
\\ 
\\ 
\texttt{Qual è il tuo prossimo comando?} \\
            }
        }
    }
    & \\ \\

    \theutterance \stepcounter{utterance}  
    & \multicolumn{4}{p{0.6\linewidth}}{
        \cellcolor[rgb]{0.9,0.9,0.9}{
            \makecell[{{p{\linewidth}}}]{
                \texttt{\tiny{[P1$\rangle$GM]}}
                \texttt{MUOVERE: A, (3,5)} \\
            }
        }
    }
    & & \\ \\

    \theutterance \stepcounter{utterance}  
    & & & \multicolumn{2}{p{0.3\linewidth}}{
        \cellcolor[rgb]{0.9,0.9,0.9}{
            \makecell[{{p{\linewidth}}}]{
                \texttt{\tiny{[GM$|$GM]}}
                \texttt{MUOVERE: A, (3,5)} \\
            }
        }
    }
    & & \\ \\

    \theutterance \stepcounter{utterance}  
    & & & \multicolumn{2}{p{0.3\linewidth}}{
        \cellcolor[rgb]{0.9,0.9,0.9}{
            \makecell[{{p{\linewidth}}}]{
                \texttt{\tiny{[GM$|$GM]}}
                \texttt{A moved to (3, 5)} \\
            }
        }
    }
    & & \\ \\

    \theutterance \stepcounter{utterance}  
    & & \multicolumn{4}{p{0.6\linewidth}}{
        \cellcolor[rgb]{0.9,0.9,0.9}{
            \makecell[{{p{\linewidth}}}]{
                \texttt{\tiny{[GM$\rangle$P2]}}
                \texttt{Il tuo messaggio è stato trasmesso all'altro giocatore.} \\
\\ 
\\ 
\texttt{Stai giocando il round 3 di un massimo di 28.} \\
\texttt{Avete accumulato collettivamente 3 di 16 penalità.} \\
\texttt{L'altro giocatore ha spostato un oggetto sulla sua griglia.} \\
\\ 
\\ 
\texttt{Qual è il tuo prossimo comando?} \\
            }
        }
    }
    & \\ \\

    \theutterance \stepcounter{utterance}  
    & & & \multicolumn{4}{p{0.6\linewidth}}{
        \cellcolor[rgb]{0.9,0.9,0.9}{
            \makecell[{{p{\linewidth}}}]{
                \texttt{\tiny{[GM$\langle$P2]}}
                \texttt{L'altro giocatore non ha ancora risposto alla mia proposta. Per ora inizio a spostare uno degli oggetti concordati. Sposto **O** in (3,2), che è vuoto e sulla quale eravamo d'accordo.} \\
\\ 
\texttt{```} \\
\texttt{MUOVERE: O, (3, 2)} \\
\texttt{```} \\
            }
        }
    }
     \\ \\

    \theutterance \stepcounter{utterance}  
    & & & \multicolumn{2}{p{0.3\linewidth}}{
        \cellcolor[rgb]{0.9,0.9,0.9}{
            \makecell[{{p{\linewidth}}}]{
                \texttt{\tiny{[GM$|$GM]}}
                \texttt{L'altro giocatore non ha ancora risposto alla mia proposta. Per ora inizio a spostare uno degli oggetti concordati. Sposto **O** in (3,2), che è vuoto e sulla quale eravamo d'accordo.} \\
\\ 
\texttt{```} \\
\texttt{MUOVERE: O, (3, 2)} \\
\texttt{```} \\
            }
        }
    }
    & & \\ \\

    \theutterance \stepcounter{utterance}  
    & & & \multicolumn{2}{p{0.3\linewidth}}{
        \cellcolor[rgb]{0.9,0.9,0.9}{
            \makecell[{{p{\linewidth}}}]{
                \texttt{\tiny{[GM$|$GM]}}
                \texttt{Invalid format: head is not empty: 'L'altro giocatore non ha ancora risposto alla mia proposta. Per ora inizio a spostare uno degli oggetti concordati. Sposto **O** in (3,2), che è vuoto e sulla quale eravamo d'accordo.} \\
\\ 
\\ 
\texttt{'} \\
            }
        }
    }
    & & \\ \\

    \theutterance \stepcounter{utterance}  
    & & \multicolumn{4}{p{0.6\linewidth}}{
        \cellcolor[rgb]{0.9,0.9,0.9}{
            \makecell[{{p{\linewidth}}}]{
                \texttt{\tiny{[GM$\rangle$P2]}}
                \texttt{Penalità: Il tuo messaggio non deve contenere nulla prima del comando!} \\
\\ 
\texttt{Assicurati che la tua risposta contenga solo `DIRE: <MESSAGGIO>` o `MUOVERE: <OGGETTO>, (<X>, <Y>)`, e nient'altro!} \\
\\ 
\texttt{Avete accumulato collettivamente 4 di 16 penalità.} \\
\texttt{Per favore riprova!} \\
            }
        }
    }
    & \\ \\

    \theutterance \stepcounter{utterance}  
    & & & \multicolumn{4}{p{0.6\linewidth}}{
        \cellcolor[rgb]{0.9,0.9,0.9}{
            \makecell[{{p{\linewidth}}}]{
                \texttt{\tiny{[GM$\langle$P2]}}
                \texttt{MUOVERE: O, (3, 2)} \\
            }
        }
    }
     \\ \\

    \theutterance \stepcounter{utterance}  
    & & & \multicolumn{2}{p{0.3\linewidth}}{
        \cellcolor[rgb]{0.9,0.9,0.9}{
            \makecell[{{p{\linewidth}}}]{
                \texttt{\tiny{[GM$|$GM]}}
                \texttt{MUOVERE: O, (3, 2)} \\
            }
        }
    }
    & & \\ \\

    \theutterance \stepcounter{utterance}  
    & & & \multicolumn{2}{p{0.3\linewidth}}{
        \cellcolor[rgb]{0.9,0.9,0.9}{
            \makecell[{{p{\linewidth}}}]{
                \texttt{\tiny{[GM$|$GM]}}
                \texttt{O moved to (3, 2)} \\
            }
        }
    }
    & & \\ \\

    \theutterance \stepcounter{utterance}  
    & & \multicolumn{4}{p{0.6\linewidth}}{
        \cellcolor[rgb]{0.9,0.9,0.9}{
            \makecell[{{p{\linewidth}}}]{
                \texttt{\tiny{[P1$\langle$GM]}}
                \texttt{Spostato 'A' in (3,5) con successo. La tua griglia aggiornata appare così:} \\
\texttt{```} \\
\texttt{ 1234567} \\
\texttt{╔═══════╗} \\
\texttt{║◌◌◌◌◌T◌║ 1} \\
\texttt{║P◌◌◌◌◌◌║ 2} \\
\texttt{║◌◌◌D◌┌─╢ 3} \\
\texttt{║◌◌◌┌─┤S║ 4} \\
\texttt{║O◌A│◌├─╢ 5} \\
\texttt{╟───┼─┘◌║ 6} \\
\texttt{║◌M◌│◌◌◌║ 7} \\
\texttt{╚═══╧═══╝} \\
\texttt{```} \\
\texttt{Stai giocando il round 3 di un massimo di 28.} \\
\texttt{Avete accumulato collettivamente 4 di 16 penalità.} \\
\texttt{L'altro giocatore ha spostato un oggetto sulla sua griglia.} \\
\\ 
\\ 
\texttt{Qual è il tuo prossimo comando?} \\
            }
        }
    }
    & \\ \\

    \theutterance \stepcounter{utterance}  
    & \multicolumn{4}{p{0.6\linewidth}}{
        \cellcolor[rgb]{0.9,0.9,0.9}{
            \makecell[{{p{\linewidth}}}]{
                \texttt{\tiny{[P1$\rangle$GM]}}
                \texttt{MUOVERE: D, (2,4)} \\
            }
        }
    }
    & & \\ \\

    \theutterance \stepcounter{utterance}  
    & & & \multicolumn{2}{p{0.3\linewidth}}{
        \cellcolor[rgb]{0.9,0.9,0.9}{
            \makecell[{{p{\linewidth}}}]{
                \texttt{\tiny{[GM$|$GM]}}
                \texttt{MUOVERE: D, (2,4)} \\
            }
        }
    }
    & & \\ \\

    \theutterance \stepcounter{utterance}  
    & & & \multicolumn{2}{p{0.3\linewidth}}{
        \cellcolor[rgb]{0.9,0.9,0.9}{
            \makecell[{{p{\linewidth}}}]{
                \texttt{\tiny{[GM$|$GM]}}
                \texttt{D moved to (2, 4)} \\
            }
        }
    }
    & & \\ \\

    \theutterance \stepcounter{utterance}  
    & & \multicolumn{4}{p{0.6\linewidth}}{
        \cellcolor[rgb]{0.9,0.9,0.9}{
            \makecell[{{p{\linewidth}}}]{
                \texttt{\tiny{[GM$\rangle$P2]}}
                \texttt{Spostato 'O' in (3,2) con successo. La tua griglia aggiornata appare così:} \\
\texttt{```} \\
\texttt{ 1234567} \\
\texttt{╔═══════╗} \\
\texttt{║◌◌◌◌◌S◌║ 1} \\
\texttt{║◌◌O◌◌◌◌║ 2} \\
\texttt{║◌◌◌M◌┌─╢ 3} \\
\texttt{║◌◌A┌─┤◌║ 4} \\
\texttt{║◌◌◌│◌├─╢ 5} \\
\texttt{╟───┼─┘◌║ 6} \\
\texttt{║D◌T│P◌◌║ 7} \\
\texttt{╚═══╧═══╝} \\
\texttt{```} \\
\texttt{Stai giocando il round 4 di un massimo di 28.} \\
\texttt{Avete accumulato collettivamente 4 di 16 penalità.} \\
\texttt{L'altro giocatore ha spostato un oggetto sulla sua griglia.} \\
\\ 
\\ 
\texttt{Qual è il tuo prossimo comando?} \\
            }
        }
    }
    & \\ \\

    \theutterance \stepcounter{utterance}  
    & & & \multicolumn{4}{p{0.6\linewidth}}{
        \cellcolor[rgb]{0.9,0.9,0.9}{
            \makecell[{{p{\linewidth}}}]{
                \texttt{\tiny{[GM$\langle$P2]}}
                \texttt{Ora sposto **D** in (2,4), che è vuoto e concordato.} \\
\\ 
\texttt{```} \\
\texttt{MUOVERE: D, (2, 4)} \\
\texttt{```} \\
            }
        }
    }
     \\ \\

    \theutterance \stepcounter{utterance}  
    & & & \multicolumn{2}{p{0.3\linewidth}}{
        \cellcolor[rgb]{0.9,0.9,0.9}{
            \makecell[{{p{\linewidth}}}]{
                \texttt{\tiny{[GM$|$GM]}}
                \texttt{Ora sposto **D** in (2,4), che è vuoto e concordato.} \\
\\ 
\texttt{```} \\
\texttt{MUOVERE: D, (2, 4)} \\
\texttt{```} \\
            }
        }
    }
    & & \\ \\

    \theutterance \stepcounter{utterance}  
    & & & \multicolumn{2}{p{0.3\linewidth}}{
        \cellcolor[rgb]{0.9,0.9,0.9}{
            \makecell[{{p{\linewidth}}}]{
                \texttt{\tiny{[GM$|$GM]}}
                \texttt{Invalid format: head is not empty: 'Ora sposto **D** in (2,4), che è vuoto e concordato.} \\
\\ 
\\ 
\texttt{'} \\
            }
        }
    }
    & & \\ \\

    \theutterance \stepcounter{utterance}  
    & & \multicolumn{4}{p{0.6\linewidth}}{
        \cellcolor[rgb]{0.9,0.9,0.9}{
            \makecell[{{p{\linewidth}}}]{
                \texttt{\tiny{[GM$\rangle$P2]}}
                \texttt{Penalità: Il tuo messaggio non deve contenere nulla prima del comando!} \\
\\ 
\texttt{Assicurati che la tua risposta contenga solo `DIRE: <MESSAGGIO>` o `MUOVERE: <OGGETTO>, (<X>, <Y>)`, e nient'altro!} \\
\\ 
\texttt{Avete accumulato collettivamente 5 di 16 penalità.} \\
\texttt{Per favore riprova!} \\
            }
        }
    }
    & \\ \\

    \theutterance \stepcounter{utterance}  
    & & & \multicolumn{4}{p{0.6\linewidth}}{
        \cellcolor[rgb]{0.9,0.9,0.9}{
            \makecell[{{p{\linewidth}}}]{
                \texttt{\tiny{[GM$\langle$P2]}}
                \texttt{MUOVERE: D, (2, 4)} \\
            }
        }
    }
     \\ \\

    \theutterance \stepcounter{utterance}  
    & & & \multicolumn{2}{p{0.3\linewidth}}{
        \cellcolor[rgb]{0.9,0.9,0.9}{
            \makecell[{{p{\linewidth}}}]{
                \texttt{\tiny{[GM$|$GM]}}
                \texttt{MUOVERE: D, (2, 4)} \\
            }
        }
    }
    & & \\ \\

    \theutterance \stepcounter{utterance}  
    & & & \multicolumn{2}{p{0.3\linewidth}}{
        \cellcolor[rgb]{0.9,0.9,0.9}{
            \makecell[{{p{\linewidth}}}]{
                \texttt{\tiny{[GM$|$GM]}}
                \texttt{D moved to (2, 4)} \\
            }
        }
    }
    & & \\ \\

    \theutterance \stepcounter{utterance}  
    & & \multicolumn{4}{p{0.6\linewidth}}{
        \cellcolor[rgb]{0.9,0.9,0.9}{
            \makecell[{{p{\linewidth}}}]{
                \texttt{\tiny{[P1$\langle$GM]}}
                \texttt{Spostato 'D' in (2,4) con successo. La tua griglia aggiornata appare così:} \\
\texttt{```} \\
\texttt{ 1234567} \\
\texttt{╔═══════╗} \\
\texttt{║◌◌◌◌◌T◌║ 1} \\
\texttt{║P◌◌◌◌◌◌║ 2} \\
\texttt{║◌◌◌◌◌┌─╢ 3} \\
\texttt{║◌D◌┌─┤S║ 4} \\
\texttt{║O◌A│◌├─╢ 5} \\
\texttt{╟───┼─┘◌║ 6} \\
\texttt{║◌M◌│◌◌◌║ 7} \\
\texttt{╚═══╧═══╝} \\
\texttt{```} \\
\texttt{Stai giocando il round 4 di un massimo di 28.} \\
\texttt{Avete accumulato collettivamente 5 di 16 penalità.} \\
\texttt{L'altro giocatore ha spostato un oggetto sulla sua griglia.} \\
\\ 
\\ 
\texttt{Qual è il tuo prossimo comando?} \\
            }
        }
    }
    & \\ \\

    \theutterance \stepcounter{utterance}  
    & \multicolumn{4}{p{0.6\linewidth}}{
        \cellcolor[rgb]{0.9,0.9,0.9}{
            \makecell[{{p{\linewidth}}}]{
                \texttt{\tiny{[P1$\rangle$GM]}}
                \texttt{MUOVERE: O, (3,2)} \\
            }
        }
    }
    & & \\ \\

    \theutterance \stepcounter{utterance}  
    & & & \multicolumn{2}{p{0.3\linewidth}}{
        \cellcolor[rgb]{0.9,0.9,0.9}{
            \makecell[{{p{\linewidth}}}]{
                \texttt{\tiny{[GM$|$GM]}}
                \texttt{MUOVERE: O, (3,2)} \\
            }
        }
    }
    & & \\ \\

    \theutterance \stepcounter{utterance}  
    & & & \multicolumn{2}{p{0.3\linewidth}}{
        \cellcolor[rgb]{0.9,0.9,0.9}{
            \makecell[{{p{\linewidth}}}]{
                \texttt{\tiny{[GM$|$GM]}}
                \texttt{O moved to (3, 2)} \\
            }
        }
    }
    & & \\ \\

    \theutterance \stepcounter{utterance}  
    & & \multicolumn{4}{p{0.6\linewidth}}{
        \cellcolor[rgb]{0.9,0.9,0.9}{
            \makecell[{{p{\linewidth}}}]{
                \texttt{\tiny{[GM$\rangle$P2]}}
                \texttt{Spostato 'D' in (2,4) con successo. La tua griglia aggiornata appare così:} \\
\texttt{```} \\
\texttt{ 1234567} \\
\texttt{╔═══════╗} \\
\texttt{║◌◌◌◌◌S◌║ 1} \\
\texttt{║◌◌O◌◌◌◌║ 2} \\
\texttt{║◌◌◌M◌┌─╢ 3} \\
\texttt{║◌DA┌─┤◌║ 4} \\
\texttt{║◌◌◌│◌├─╢ 5} \\
\texttt{╟───┼─┘◌║ 6} \\
\texttt{║◌◌T│P◌◌║ 7} \\
\texttt{╚═══╧═══╝} \\
\texttt{```} \\
\texttt{Stai giocando il round 5 di un massimo di 28.} \\
\texttt{Avete accumulato collettivamente 5 di 16 penalità.} \\
\texttt{L'altro giocatore ha spostato un oggetto sulla sua griglia.} \\
\\ 
\\ 
\texttt{Qual è il tuo prossimo comando?} \\
            }
        }
    }
    & \\ \\

    \theutterance \stepcounter{utterance}  
    & & & \multicolumn{4}{p{0.6\linewidth}}{
        \cellcolor[rgb]{0.9,0.9,0.9}{
            \makecell[{{p{\linewidth}}}]{
                \texttt{\tiny{[GM$\langle$P2]}}
                \texttt{Ora sposto **T** in (6,4), che è vuoto e concordato.} \\
\\ 
\texttt{MUOVERE: T, (6, 4)} \\
            }
        }
    }
     \\ \\

    \theutterance \stepcounter{utterance}  
    & & & \multicolumn{2}{p{0.3\linewidth}}{
        \cellcolor[rgb]{0.9,0.9,0.9}{
            \makecell[{{p{\linewidth}}}]{
                \texttt{\tiny{[GM$|$GM]}}
                \texttt{Ora sposto **T** in (6,4), che è vuoto e concordato.} \\
\\ 
\texttt{MUOVERE: T, (6, 4)} \\
            }
        }
    }
    & & \\ \\

    \theutterance \stepcounter{utterance}  
    & & & \multicolumn{2}{p{0.3\linewidth}}{
        \cellcolor[rgb]{0.9,0.9,0.9}{
            \makecell[{{p{\linewidth}}}]{
                \texttt{\tiny{[GM$|$GM]}}
                \texttt{Invalid format: head is not empty: 'Ora sposto **T** in (6,4), che è vuoto e concordato.} \\
\\ 
\texttt{'} \\
            }
        }
    }
    & & \\ \\

    \theutterance \stepcounter{utterance}  
    & & \multicolumn{4}{p{0.6\linewidth}}{
        \cellcolor[rgb]{0.9,0.9,0.9}{
            \makecell[{{p{\linewidth}}}]{
                \texttt{\tiny{[GM$\rangle$P2]}}
                \texttt{Penalità: Il tuo messaggio non deve contenere nulla prima del comando!} \\
\\ 
\texttt{Assicurati che la tua risposta contenga solo `DIRE: <MESSAGGIO>` o `MUOVERE: <OGGETTO>, (<X>, <Y>)`, e nient'altro!} \\
\\ 
\texttt{Avete accumulato collettivamente 6 di 16 penalità.} \\
\texttt{Per favore riprova!} \\
            }
        }
    }
    & \\ \\

    \theutterance \stepcounter{utterance}  
    & & & \multicolumn{4}{p{0.6\linewidth}}{
        \cellcolor[rgb]{0.9,0.9,0.9}{
            \makecell[{{p{\linewidth}}}]{
                \texttt{\tiny{[GM$\langle$P2]}}
                \texttt{MUOVERE: T, (6, 4)} \\
            }
        }
    }
     \\ \\

    \theutterance \stepcounter{utterance}  
    & & & \multicolumn{2}{p{0.3\linewidth}}{
        \cellcolor[rgb]{0.9,0.9,0.9}{
            \makecell[{{p{\linewidth}}}]{
                \texttt{\tiny{[GM$|$GM]}}
                \texttt{MUOVERE: T, (6, 4)} \\
            }
        }
    }
    & & \\ \\

    \theutterance \stepcounter{utterance}  
    & & & \multicolumn{2}{p{0.3\linewidth}}{
        \cellcolor[rgb]{0.9,0.9,0.9}{
            \makecell[{{p{\linewidth}}}]{
                \texttt{\tiny{[GM$|$GM]}}
                \texttt{Penalità: (6,4) non è vuoto, ma contiene '┤'.} \\
\\ 
\texttt{Avete accumulato collettivamente 7 di 16 penalità.} \\
\texttt{Per favore riprova!} \\
            }
        }
    }
    & & \\ \\

    \theutterance \stepcounter{utterance}  
    & & \multicolumn{4}{p{0.6\linewidth}}{
        \cellcolor[rgb]{0.9,0.9,0.9}{
            \makecell[{{p{\linewidth}}}]{
                \texttt{\tiny{[GM$\rangle$P2]}}
                \texttt{Penalità: (6,4) non è vuoto, ma contiene '┤'.} \\
\\ 
\texttt{Avete accumulato collettivamente 7 di 16 penalità.} \\
\texttt{Per favore riprova!} \\
            }
        }
    }
    & \\ \\

    \theutterance \stepcounter{utterance}  
    & & & \multicolumn{4}{p{0.6\linewidth}}{
        \cellcolor[rgb]{0.9,0.9,0.9}{
            \makecell[{{p{\linewidth}}}]{
                \texttt{\tiny{[GM$\langle$P2]}}
                \texttt{Prima di spostare T, devo controllare la griglia attuale. Alla riga 4, colonna 6 c'è `┤`. Devo comunicare all'altro giocatore che (6,4) non è vuoto e proporre un'alternativa, come (7,4) che è vuoto.} \\
\\ 
\texttt{DIRE: (6,4) da me non è vuoto, contiene una linea. Propongo T a (7,4) che è libero. Confermi?} \\
            }
        }
    }
     \\ \\

    \theutterance \stepcounter{utterance}  
    & & & \multicolumn{2}{p{0.3\linewidth}}{
        \cellcolor[rgb]{0.9,0.9,0.9}{
            \makecell[{{p{\linewidth}}}]{
                \texttt{\tiny{[GM$|$GM]}}
                \texttt{Prima di spostare T, devo controllare la griglia attuale. Alla riga 4, colonna 6 c'è `┤`. Devo comunicare all'altro giocatore che (6,4) non è vuoto e proporre un'alternativa, come (7,4) che è vuoto.} \\
\\ 
\texttt{DIRE: (6,4) da me non è vuoto, contiene una linea. Propongo T a (7,4) che è libero. Confermi?} \\
            }
        }
    }
    & & \\ \\

    \theutterance \stepcounter{utterance}  
    & & & \multicolumn{2}{p{0.3\linewidth}}{
        \cellcolor[rgb]{0.9,0.9,0.9}{
            \makecell[{{p{\linewidth}}}]{
                \texttt{\tiny{[GM$|$GM]}}
                \texttt{Invalid format: head is not empty: 'Prima di spostare T, devo controllare la griglia attuale. Alla riga 4, colonna 6 c'è ┤. Devo comunicare all'altro giocatore che (6,4) non è vuoto e proporre un'alternativa, come (7,4) che è vuoto.} \\
\\ 
\texttt{'} \\
            }
        }
    }
    & & \\ \\

    \theutterance \stepcounter{utterance}  
    & & \multicolumn{4}{p{0.6\linewidth}}{
        \cellcolor[rgb]{0.9,0.9,0.9}{
            \makecell[{{p{\linewidth}}}]{
                \texttt{\tiny{[GM$\rangle$P2]}}
                \texttt{Penalità: Il tuo messaggio non deve contenere nulla prima del comando!} \\
\\ 
\texttt{Assicurati che la tua risposta contenga solo `DIRE: <MESSAGGIO>` o `MUOVERE: <OGGETTO>, (<X>, <Y>)`, e nient'altro!} \\
\\ 
\texttt{Avete accumulato collettivamente 8 di 16 penalità.} \\
\texttt{Per favore riprova!} \\
            }
        }
    }
    & \\ \\

    \theutterance \stepcounter{utterance}  
    & & & \multicolumn{4}{p{0.6\linewidth}}{
        \cellcolor[rgb]{0.9,0.9,0.9}{
            \makecell[{{p{\linewidth}}}]{
                \texttt{\tiny{[GM$\langle$P2]}}
                \texttt{DIRE: (6,4) da me non è vuoto, contiene una linea. Propongo T a (7,4) che è libero. Confermi?} \\
            }
        }
    }
     \\ \\

    \theutterance \stepcounter{utterance}  
    & & & \multicolumn{2}{p{0.3\linewidth}}{
        \cellcolor[rgb]{0.9,0.9,0.9}{
            \makecell[{{p{\linewidth}}}]{
                \texttt{\tiny{[GM$|$GM]}}
                \texttt{DIRE: (6,4) da me non è vuoto, contiene una linea. Propongo T a (7,4) che è libero. Confermi?} \\
            }
        }
    }
    & & \\ \\

    \theutterance \stepcounter{utterance}  
    & & \multicolumn{4}{p{0.6\linewidth}}{
        \cellcolor[rgb]{0.9,0.9,0.9}{
            \makecell[{{p{\linewidth}}}]{
                \texttt{\tiny{[P1$\langle$GM]}}
                \texttt{Spostato 'O' in (3,2) con successo. La tua griglia aggiornata appare così:} \\
\texttt{```} \\
\texttt{ 1234567} \\
\texttt{╔═══════╗} \\
\texttt{║◌◌◌◌◌T◌║ 1} \\
\texttt{║P◌O◌◌◌◌║ 2} \\
\texttt{║◌◌◌◌◌┌─╢ 3} \\
\texttt{║◌D◌┌─┤S║ 4} \\
\texttt{║◌◌A│◌├─╢ 5} \\
\texttt{╟───┼─┘◌║ 6} \\
\texttt{║◌M◌│◌◌◌║ 7} \\
\texttt{╚═══╧═══╝} \\
\texttt{```} \\
\texttt{Stai giocando il round 5 di un massimo di 28.} \\
\texttt{Avete accumulato collettivamente 8 di 16 penalità.} \\
\texttt{L'altro giocatore ha inviato questo messaggio:} \\
\\ 
\\ 
\texttt{"(6,4) da me non è vuoto, contiene una linea. Propongo T a (7,4) che è libero. Confermi?"} \\
\\ 
\\ 
\texttt{Qual è il tuo prossimo comando?} \\
            }
        }
    }
    & \\ \\

    \theutterance \stepcounter{utterance}  
    & \multicolumn{4}{p{0.6\linewidth}}{
        \cellcolor[rgb]{0.9,0.9,0.9}{
            \makecell[{{p{\linewidth}}}]{
                \texttt{\tiny{[P1$\rangle$GM]}}
                \texttt{DIRE: Sì, T a (7,4) va bene per me. Confermo.} \\
            }
        }
    }
    & & \\ \\

    \theutterance \stepcounter{utterance}  
    & & & \multicolumn{2}{p{0.3\linewidth}}{
        \cellcolor[rgb]{0.9,0.9,0.9}{
            \makecell[{{p{\linewidth}}}]{
                \texttt{\tiny{[GM$|$GM]}}
                \texttt{DIRE: Sì, T a (7,4) va bene per me. Confermo.} \\
            }
        }
    }
    & & \\ \\

    \theutterance \stepcounter{utterance}  
    & & \multicolumn{4}{p{0.6\linewidth}}{
        \cellcolor[rgb]{0.9,0.9,0.9}{
            \makecell[{{p{\linewidth}}}]{
                \texttt{\tiny{[GM$\rangle$P2]}}
                \texttt{Il tuo messaggio è stato trasmesso all'altro giocatore.} \\
\\ 
\\ 
\texttt{Stai giocando il round 6 di un massimo di 28.} \\
\texttt{Avete accumulato collettivamente 8 di 16 penalità.} \\
\texttt{L'altro giocatore ha inviato questo messaggio:} \\
\\ 
\\ 
\texttt{"Sì, T a (7,4) va bene per me. Confermo."} \\
\\ 
\\ 
\texttt{Qual è il tuo prossimo comando?} \\
            }
        }
    }
    & \\ \\

    \theutterance \stepcounter{utterance}  
    & & & \multicolumn{4}{p{0.6\linewidth}}{
        \cellcolor[rgb]{0.9,0.9,0.9}{
            \makecell[{{p{\linewidth}}}]{
                \texttt{\tiny{[GM$\langle$P2]}}
                \texttt{MUOVERE: T, (7, 4)} \\
            }
        }
    }
     \\ \\

    \theutterance \stepcounter{utterance}  
    & & & \multicolumn{2}{p{0.3\linewidth}}{
        \cellcolor[rgb]{0.9,0.9,0.9}{
            \makecell[{{p{\linewidth}}}]{
                \texttt{\tiny{[GM$|$GM]}}
                \texttt{MUOVERE: T, (7, 4)} \\
            }
        }
    }
    & & \\ \\

    \theutterance \stepcounter{utterance}  
    & & & \multicolumn{2}{p{0.3\linewidth}}{
        \cellcolor[rgb]{0.9,0.9,0.9}{
            \makecell[{{p{\linewidth}}}]{
                \texttt{\tiny{[GM$|$GM]}}
                \texttt{T moved to (7, 4)} \\
            }
        }
    }
    & & \\ \\

    \theutterance \stepcounter{utterance}  
    & & \multicolumn{4}{p{0.6\linewidth}}{
        \cellcolor[rgb]{0.9,0.9,0.9}{
            \makecell[{{p{\linewidth}}}]{
                \texttt{\tiny{[P1$\langle$GM]}}
                \texttt{Il tuo messaggio è stato trasmesso all'altro giocatore.} \\
\\ 
\\ 
\texttt{Stai giocando il round 6 di un massimo di 28.} \\
\texttt{Avete accumulato collettivamente 8 di 16 penalità.} \\
\texttt{L'altro giocatore ha spostato un oggetto sulla sua griglia.} \\
\\ 
\\ 
\texttt{Qual è il tuo prossimo comando?} \\
            }
        }
    }
    & \\ \\

    \theutterance \stepcounter{utterance}  
    & \multicolumn{4}{p{0.6\linewidth}}{
        \cellcolor[rgb]{0.9,0.9,0.9}{
            \makecell[{{p{\linewidth}}}]{
                \texttt{\tiny{[P1$\rangle$GM]}}
                \texttt{MUOVERE: T, (7,4)} \\
            }
        }
    }
    & & \\ \\

    \theutterance \stepcounter{utterance}  
    & & & \multicolumn{2}{p{0.3\linewidth}}{
        \cellcolor[rgb]{0.9,0.9,0.9}{
            \makecell[{{p{\linewidth}}}]{
                \texttt{\tiny{[GM$|$GM]}}
                \texttt{MUOVERE: T, (7,4)} \\
            }
        }
    }
    & & \\ \\

    \theutterance \stepcounter{utterance}  
    & & & \multicolumn{2}{p{0.3\linewidth}}{
        \cellcolor[rgb]{0.9,0.9,0.9}{
            \makecell[{{p{\linewidth}}}]{
                \texttt{\tiny{[GM$|$GM]}}
                \texttt{Penalità: (7,4) non è vuoto, ma contiene 'S'.} \\
\\ 
\texttt{Avete accumulato collettivamente 9 di 16 penalità.} \\
\texttt{Per favore riprova!} \\
            }
        }
    }
    & & \\ \\

    \theutterance \stepcounter{utterance}  
    & & \multicolumn{4}{p{0.6\linewidth}}{
        \cellcolor[rgb]{0.9,0.9,0.9}{
            \makecell[{{p{\linewidth}}}]{
                \texttt{\tiny{[P1$\langle$GM]}}
                \texttt{Penalità: (7,4) non è vuoto, ma contiene 'S'.} \\
\\ 
\texttt{Avete accumulato collettivamente 9 di 16 penalità.} \\
\texttt{Per favore riprova!} \\
            }
        }
    }
    & \\ \\

    \theutterance \stepcounter{utterance}  
    & \multicolumn{4}{p{0.6\linewidth}}{
        \cellcolor[rgb]{0.9,0.9,0.9}{
            \makecell[{{p{\linewidth}}}]{
                \texttt{\tiny{[P1$\rangle$GM]}}
                \texttt{MUOVERE: T, (7,3)} \\
            }
        }
    }
    & & \\ \\

    \theutterance \stepcounter{utterance}  
    & & & \multicolumn{2}{p{0.3\linewidth}}{
        \cellcolor[rgb]{0.9,0.9,0.9}{
            \makecell[{{p{\linewidth}}}]{
                \texttt{\tiny{[GM$|$GM]}}
                \texttt{MUOVERE: T, (7,3)} \\
            }
        }
    }
    & & \\ \\

    \theutterance \stepcounter{utterance}  
    & & & \multicolumn{2}{p{0.3\linewidth}}{
        \cellcolor[rgb]{0.9,0.9,0.9}{
            \makecell[{{p{\linewidth}}}]{
                \texttt{\tiny{[GM$|$GM]}}
                \texttt{Penalità: (7,3) non è vuoto, ma contiene '─'.} \\
\\ 
\texttt{Avete accumulato collettivamente 10 di 16 penalità.} \\
\texttt{Per favore riprova!} \\
            }
        }
    }
    & & \\ \\

    \theutterance \stepcounter{utterance}  
    & & \multicolumn{4}{p{0.6\linewidth}}{
        \cellcolor[rgb]{0.9,0.9,0.9}{
            \makecell[{{p{\linewidth}}}]{
                \texttt{\tiny{[P1$\langle$GM]}}
                \texttt{Penalità: (7,3) non è vuoto, ma contiene '─'.} \\
\\ 
\texttt{Avete accumulato collettivamente 10 di 16 penalità.} \\
\texttt{Per favore riprova!} \\
            }
        }
    }
    & \\ \\

    \theutterance \stepcounter{utterance}  
    & \multicolumn{4}{p{0.6\linewidth}}{
        \cellcolor[rgb]{0.9,0.9,0.9}{
            \makecell[{{p{\linewidth}}}]{
                \texttt{\tiny{[P1$\rangle$GM]}}
                \texttt{MUOVERE: T, (5,4)} \\
            }
        }
    }
    & & \\ \\

    \theutterance \stepcounter{utterance}  
    & & & \multicolumn{2}{p{0.3\linewidth}}{
        \cellcolor[rgb]{0.9,0.9,0.9}{
            \makecell[{{p{\linewidth}}}]{
                \texttt{\tiny{[GM$|$GM]}}
                \texttt{MUOVERE: T, (5,4)} \\
            }
        }
    }
    & & \\ \\

    \theutterance \stepcounter{utterance}  
    & & & \multicolumn{2}{p{0.3\linewidth}}{
        \cellcolor[rgb]{0.9,0.9,0.9}{
            \makecell[{{p{\linewidth}}}]{
                \texttt{\tiny{[GM$|$GM]}}
                \texttt{Penalità: (5,4) non è vuoto, ma contiene '─'.} \\
\\ 
\texttt{Avete accumulato collettivamente 11 di 16 penalità.} \\
\texttt{Per favore riprova!} \\
            }
        }
    }
    & & \\ \\

    \theutterance \stepcounter{utterance}  
    & & \multicolumn{4}{p{0.6\linewidth}}{
        \cellcolor[rgb]{0.9,0.9,0.9}{
            \makecell[{{p{\linewidth}}}]{
                \texttt{\tiny{[P1$\langle$GM]}}
                \texttt{Penalità: (5,4) non è vuoto, ma contiene '─'.} \\
\\ 
\texttt{Avete accumulato collettivamente 11 di 16 penalità.} \\
\texttt{Per favore riprova!} \\
            }
        }
    }
    & \\ \\

    \theutterance \stepcounter{utterance}  
    & \multicolumn{4}{p{0.6\linewidth}}{
        \cellcolor[rgb]{0.9,0.9,0.9}{
            \makecell[{{p{\linewidth}}}]{
                \texttt{\tiny{[P1$\rangle$GM]}}
                \texttt{MUOVERE: T, (7,1)} \\
            }
        }
    }
    & & \\ \\

    \theutterance \stepcounter{utterance}  
    & & & \multicolumn{2}{p{0.3\linewidth}}{
        \cellcolor[rgb]{0.9,0.9,0.9}{
            \makecell[{{p{\linewidth}}}]{
                \texttt{\tiny{[GM$|$GM]}}
                \texttt{MUOVERE: T, (7,1)} \\
            }
        }
    }
    & & \\ \\

    \theutterance \stepcounter{utterance}  
    & & & \multicolumn{2}{p{0.3\linewidth}}{
        \cellcolor[rgb]{0.9,0.9,0.9}{
            \makecell[{{p{\linewidth}}}]{
                \texttt{\tiny{[GM$|$GM]}}
                \texttt{T moved to (7, 1)} \\
            }
        }
    }
    & & \\ \\

    \theutterance \stepcounter{utterance}  
    & & \multicolumn{4}{p{0.6\linewidth}}{
        \cellcolor[rgb]{0.9,0.9,0.9}{
            \makecell[{{p{\linewidth}}}]{
                \texttt{\tiny{[GM$\rangle$P2]}}
                \texttt{Spostato 'T' in (7,4) con successo. La tua griglia aggiornata appare così:} \\
\texttt{```} \\
\texttt{ 1234567} \\
\texttt{╔═══════╗} \\
\texttt{║◌◌◌◌◌S◌║ 1} \\
\texttt{║◌◌O◌◌◌◌║ 2} \\
\texttt{║◌◌◌M◌┌─╢ 3} \\
\texttt{║◌DA┌─┤T║ 4} \\
\texttt{║◌◌◌│◌├─╢ 5} \\
\texttt{╟───┼─┘◌║ 6} \\
\texttt{║◌◌◌│P◌◌║ 7} \\
\texttt{╚═══╧═══╝} \\
\texttt{```} \\
\texttt{Stai giocando il round 7 di un massimo di 28.} \\
\texttt{Avete accumulato collettivamente 11 di 16 penalità.} \\
\texttt{L'altro giocatore ha spostato un oggetto sulla sua griglia.} \\
\\ 
\\ 
\texttt{Qual è il tuo prossimo comando?} \\
            }
        }
    }
    & \\ \\

    \theutterance \stepcounter{utterance}  
    & & & \multicolumn{4}{p{0.6\linewidth}}{
        \cellcolor[rgb]{0.9,0.9,0.9}{
            \makecell[{{p{\linewidth}}}]{
                \texttt{\tiny{[GM$\langle$P2]}}
                \texttt{Ora sposto **A** in (3,5). Controllo la griglia: riga 5, colonna 3 è `◌` – vuoto.} \\
\\ 
\texttt{MUOVERE: A, (3, 5)} \\
            }
        }
    }
     \\ \\

    \theutterance \stepcounter{utterance}  
    & & & \multicolumn{2}{p{0.3\linewidth}}{
        \cellcolor[rgb]{0.9,0.9,0.9}{
            \makecell[{{p{\linewidth}}}]{
                \texttt{\tiny{[GM$|$GM]}}
                \texttt{Ora sposto **A** in (3,5). Controllo la griglia: riga 5, colonna 3 è `◌` – vuoto.} \\
\\ 
\texttt{MUOVERE: A, (3, 5)} \\
            }
        }
    }
    & & \\ \\

    \theutterance \stepcounter{utterance}  
    & & & \multicolumn{2}{p{0.3\linewidth}}{
        \cellcolor[rgb]{0.9,0.9,0.9}{
            \makecell[{{p{\linewidth}}}]{
                \texttt{\tiny{[GM$|$GM]}}
                \texttt{Invalid format: head is not empty: 'Ora sposto **A** in (3,5). Controllo la griglia: riga 5, colonna 3 è ◌ – vuoto.} \\
\\ 
\texttt{'} \\
            }
        }
    }
    & & \\ \\

    \theutterance \stepcounter{utterance}  
    & & \multicolumn{4}{p{0.6\linewidth}}{
        \cellcolor[rgb]{0.9,0.9,0.9}{
            \makecell[{{p{\linewidth}}}]{
                \texttt{\tiny{[GM$\rangle$P2]}}
                \texttt{Penalità: Il tuo messaggio non deve contenere nulla prima del comando!} \\
\\ 
\texttt{Assicurati che la tua risposta contenga solo `DIRE: <MESSAGGIO>` o `MUOVERE: <OGGETTO>, (<X>, <Y>)`, e nient'altro!} \\
\\ 
\texttt{Avete accumulato collettivamente 12 di 16 penalità.} \\
\texttt{Per favore riprova!} \\
            }
        }
    }
    & \\ \\

    \theutterance \stepcounter{utterance}  
    & & & \multicolumn{4}{p{0.6\linewidth}}{
        \cellcolor[rgb]{0.9,0.9,0.9}{
            \makecell[{{p{\linewidth}}}]{
                \texttt{\tiny{[GM$\langle$P2]}}
                \texttt{MUOVERE: A, (3, 5)} \\
            }
        }
    }
     \\ \\

    \theutterance \stepcounter{utterance}  
    & & & \multicolumn{2}{p{0.3\linewidth}}{
        \cellcolor[rgb]{0.9,0.9,0.9}{
            \makecell[{{p{\linewidth}}}]{
                \texttt{\tiny{[GM$|$GM]}}
                \texttt{MUOVERE: A, (3, 5)} \\
            }
        }
    }
    & & \\ \\

    \theutterance \stepcounter{utterance}  
    & & & \multicolumn{2}{p{0.3\linewidth}}{
        \cellcolor[rgb]{0.9,0.9,0.9}{
            \makecell[{{p{\linewidth}}}]{
                \texttt{\tiny{[GM$|$GM]}}
                \texttt{A moved to (3, 5)} \\
            }
        }
    }
    & & \\ \\

    \theutterance \stepcounter{utterance}  
    & & \multicolumn{4}{p{0.6\linewidth}}{
        \cellcolor[rgb]{0.9,0.9,0.9}{
            \makecell[{{p{\linewidth}}}]{
                \texttt{\tiny{[P1$\langle$GM]}}
                \texttt{Spostato 'T' in (7,1) con successo. La tua griglia aggiornata appare così:} \\
\texttt{```} \\
\texttt{ 1234567} \\
\texttt{╔═══════╗} \\
\texttt{║◌◌◌◌◌◌T║ 1} \\
\texttt{║P◌O◌◌◌◌║ 2} \\
\texttt{║◌◌◌◌◌┌─╢ 3} \\
\texttt{║◌D◌┌─┤S║ 4} \\
\texttt{║◌◌A│◌├─╢ 5} \\
\texttt{╟───┼─┘◌║ 6} \\
\texttt{║◌M◌│◌◌◌║ 7} \\
\texttt{╚═══╧═══╝} \\
\texttt{```} \\
\texttt{Stai giocando il round 7 di un massimo di 28.} \\
\texttt{Avete accumulato collettivamente 12 di 16 penalità.} \\
\texttt{L'altro giocatore ha spostato un oggetto sulla sua griglia.} \\
\\ 
\\ 
\texttt{Qual è il tuo prossimo comando?} \\
            }
        }
    }
    & \\ \\

    \theutterance \stepcounter{utterance}  
    & \multicolumn{4}{p{0.6\linewidth}}{
        \cellcolor[rgb]{0.9,0.9,0.9}{
            \makecell[{{p{\linewidth}}}]{
                \texttt{\tiny{[P1$\rangle$GM]}}
                \texttt{DIRE: T l'ho spostato a (7,1) perché (7,4) e (7,3) erano occupati. Per T dove lo mettiamo? Propongo di tenerlo a (7,1) se per te va bene.} \\
            }
        }
    }
    & & \\ \\

    \theutterance \stepcounter{utterance}  
    & & & \multicolumn{2}{p{0.3\linewidth}}{
        \cellcolor[rgb]{0.9,0.9,0.9}{
            \makecell[{{p{\linewidth}}}]{
                \texttt{\tiny{[GM$|$GM]}}
                \texttt{DIRE: T l'ho spostato a (7,1) perché (7,4) e (7,3) erano occupati. Per T dove lo mettiamo? Propongo di tenerlo a (7,1) se per te va bene.} \\
            }
        }
    }
    & & \\ \\

    \theutterance \stepcounter{utterance}  
    & & \multicolumn{4}{p{0.6\linewidth}}{
        \cellcolor[rgb]{0.9,0.9,0.9}{
            \makecell[{{p{\linewidth}}}]{
                \texttt{\tiny{[GM$\rangle$P2]}}
                \texttt{Spostato 'A' in (3,5) con successo. La tua griglia aggiornata appare così:} \\
\texttt{```} \\
\texttt{ 1234567} \\
\texttt{╔═══════╗} \\
\texttt{║◌◌◌◌◌S◌║ 1} \\
\texttt{║◌◌O◌◌◌◌║ 2} \\
\texttt{║◌◌◌M◌┌─╢ 3} \\
\texttt{║◌D◌┌─┤T║ 4} \\
\texttt{║◌◌A│◌├─╢ 5} \\
\texttt{╟───┼─┘◌║ 6} \\
\texttt{║◌◌◌│P◌◌║ 7} \\
\texttt{╚═══╧═══╝} \\
\texttt{```} \\
\texttt{Stai giocando il round 8 di un massimo di 28.} \\
\texttt{Avete accumulato collettivamente 12 di 16 penalità.} \\
\texttt{L'altro giocatore ha inviato questo messaggio:} \\
\\ 
\\ 
\texttt{"T l'ho spostato a (7,1) perché (7,4) e (7,3) erano occupati. Per T dove lo mettiamo? Propongo di tenerlo a (7,1) se per te va bene."} \\
\\ 
\\ 
\texttt{Qual è il tuo prossimo comando?} \\
            }
        }
    }
    & \\ \\

    \theutterance \stepcounter{utterance}  
    & & & \multicolumn{4}{p{0.6\linewidth}}{
        \cellcolor[rgb]{0.9,0.9,0.9}{
            \makecell[{{p{\linewidth}}}]{
                \texttt{\tiny{[GM$\langle$P2]}}
                \texttt{L'altro giocatore ha spostato T a (7,1). Sulla mia griglia (7,1) è vuoto (`◌`). Per me va bene. Gli rispondo e intanto sposto **M** in (3,6) – ma devo controllare: riga 6, colonna 3 ha `─`, quindi non è vuoto. Propongo alternativa.} \\
\\ 
\texttt{DIRE: T a (7,1) per me va bene. Per M, (3,6) da me ha una linea. Propongo M a (2,6) che è vuoto. Confermi?} \\
            }
        }
    }
     \\ \\

    \theutterance \stepcounter{utterance}  
    & & & \multicolumn{2}{p{0.3\linewidth}}{
        \cellcolor[rgb]{0.9,0.9,0.9}{
            \makecell[{{p{\linewidth}}}]{
                \texttt{\tiny{[GM$|$GM]}}
                \texttt{L'altro giocatore ha spostato T a (7,1). Sulla mia griglia (7,1) è vuoto (`◌`). Per me va bene. Gli rispondo e intanto sposto **M** in (3,6) – ma devo controllare: riga 6, colonna 3 ha `─`, quindi non è vuoto. Propongo alternativa.} \\
\\ 
\texttt{DIRE: T a (7,1) per me va bene. Per M, (3,6) da me ha una linea. Propongo M a (2,6) che è vuoto. Confermi?} \\
            }
        }
    }
    & & \\ \\

    \theutterance \stepcounter{utterance}  
    & & & \multicolumn{2}{p{0.3\linewidth}}{
        \cellcolor[rgb]{0.9,0.9,0.9}{
            \makecell[{{p{\linewidth}}}]{
                \texttt{\tiny{[GM$|$GM]}}
                \texttt{Invalid format: head is not empty: 'L'altro giocatore ha spostato T a (7,1). Sulla mia griglia (7,1) è vuoto (◌). Per me va bene. Gli rispondo e intanto sposto **M** in (3,6) – ma devo controllare: riga 6, colonna 3 ha ─, quindi non è vuoto. Propongo alternativa.} \\
\\ 
\texttt{'} \\
            }
        }
    }
    & & \\ \\

    \theutterance \stepcounter{utterance}  
    & & \multicolumn{4}{p{0.6\linewidth}}{
        \cellcolor[rgb]{0.9,0.9,0.9}{
            \makecell[{{p{\linewidth}}}]{
                \texttt{\tiny{[GM$\rangle$P2]}}
                \texttt{Penalità: Il tuo messaggio non deve contenere nulla prima del comando!} \\
\\ 
\texttt{Assicurati che la tua risposta contenga solo `DIRE: <MESSAGGIO>` o `MUOVERE: <OGGETTO>, (<X>, <Y>)`, e nient'altro!} \\
\\ 
\texttt{Avete accumulato collettivamente 13 di 16 penalità.} \\
\texttt{Per favore riprova!} \\
            }
        }
    }
    & \\ \\

    \theutterance \stepcounter{utterance}  
    & & & \multicolumn{4}{p{0.6\linewidth}}{
        \cellcolor[rgb]{0.9,0.9,0.9}{
            \makecell[{{p{\linewidth}}}]{
                \texttt{\tiny{[GM$\langle$P2]}}
                \texttt{DIRE: T a (7,1) per me va bene. Per M, (3,6) da me ha una linea. Propongo M a (2,6) che è vuoto. Confermi?} \\
            }
        }
    }
     \\ \\

    \theutterance \stepcounter{utterance}  
    & & & \multicolumn{2}{p{0.3\linewidth}}{
        \cellcolor[rgb]{0.9,0.9,0.9}{
            \makecell[{{p{\linewidth}}}]{
                \texttt{\tiny{[GM$|$GM]}}
                \texttt{DIRE: T a (7,1) per me va bene. Per M, (3,6) da me ha una linea. Propongo M a (2,6) che è vuoto. Confermi?} \\
            }
        }
    }
    & & \\ \\

    \theutterance \stepcounter{utterance}  
    & & \multicolumn{4}{p{0.6\linewidth}}{
        \cellcolor[rgb]{0.9,0.9,0.9}{
            \makecell[{{p{\linewidth}}}]{
                \texttt{\tiny{[P1$\langle$GM]}}
                \texttt{Il tuo messaggio è stato trasmesso all'altro giocatore.} \\
\\ 
\\ 
\texttt{Stai giocando il round 8 di un massimo di 28.} \\
\texttt{Avete accumulato collettivamente 13 di 16 penalità.} \\
\texttt{L'altro giocatore ha inviato questo messaggio:} \\
\\ 
\\ 
\texttt{"T a (7,1) per me va bene. Per M, (3,6) da me ha una linea. Propongo M a (2,6) che è vuoto. Confermi?"} \\
\\ 
\\ 
\texttt{Qual è il tuo prossimo comando?} \\
            }
        }
    }
    & \\ \\

    \theutterance \stepcounter{utterance}  
    & \multicolumn{4}{p{0.6\linewidth}}{
        \cellcolor[rgb]{0.9,0.9,0.9}{
            \makecell[{{p{\linewidth}}}]{
                \texttt{\tiny{[P1$\rangle$GM]}}
                \texttt{MUOVERE: M, (2,6)} \\
            }
        }
    }
    & & \\ \\

    \theutterance \stepcounter{utterance}  
    & & & \multicolumn{2}{p{0.3\linewidth}}{
        \cellcolor[rgb]{0.9,0.9,0.9}{
            \makecell[{{p{\linewidth}}}]{
                \texttt{\tiny{[GM$|$GM]}}
                \texttt{MUOVERE: M, (2,6)} \\
            }
        }
    }
    & & \\ \\

    \theutterance \stepcounter{utterance}  
    & & & \multicolumn{2}{p{0.3\linewidth}}{
        \cellcolor[rgb]{0.9,0.9,0.9}{
            \makecell[{{p{\linewidth}}}]{
                \texttt{\tiny{[GM$|$GM]}}
                \texttt{Penalità: (2,6) non è vuoto, ma contiene '─'.} \\
\\ 
\texttt{Avete accumulato collettivamente 14 di 16 penalità.} \\
\texttt{Per favore riprova!} \\
            }
        }
    }
    & & \\ \\

    \theutterance \stepcounter{utterance}  
    & & \multicolumn{4}{p{0.6\linewidth}}{
        \cellcolor[rgb]{0.9,0.9,0.9}{
            \makecell[{{p{\linewidth}}}]{
                \texttt{\tiny{[P1$\langle$GM]}}
                \texttt{Penalità: (2,6) non è vuoto, ma contiene '─'.} \\
\\ 
\texttt{Avete accumulato collettivamente 14 di 16 penalità.} \\
\texttt{Per favore riprova!} \\
            }
        }
    }
    & \\ \\

    \theutterance \stepcounter{utterance}  
    & \multicolumn{4}{p{0.6\linewidth}}{
        \cellcolor[rgb]{0.9,0.9,0.9}{
            \makecell[{{p{\linewidth}}}]{
                \texttt{\tiny{[P1$\rangle$GM]}}
                \texttt{MUOVERE: M, (1,6)} \\
            }
        }
    }
    & & \\ \\

    \theutterance \stepcounter{utterance}  
    & & & \multicolumn{2}{p{0.3\linewidth}}{
        \cellcolor[rgb]{0.9,0.9,0.9}{
            \makecell[{{p{\linewidth}}}]{
                \texttt{\tiny{[GM$|$GM]}}
                \texttt{MUOVERE: M, (1,6)} \\
            }
        }
    }
    & & \\ \\

    \theutterance \stepcounter{utterance}  
    & & & \multicolumn{2}{p{0.3\linewidth}}{
        \cellcolor[rgb]{0.9,0.9,0.9}{
            \makecell[{{p{\linewidth}}}]{
                \texttt{\tiny{[GM$|$GM]}}
                \texttt{Penalità: (1,6) non è vuoto, ma contiene '─'.} \\
\\ 
\texttt{Avete accumulato collettivamente 15 di 16 penalità.} \\
\texttt{Per favore riprova!} \\
            }
        }
    }
    & & \\ \\

    \theutterance \stepcounter{utterance}  
    & & \multicolumn{4}{p{0.6\linewidth}}{
        \cellcolor[rgb]{0.9,0.9,0.9}{
            \makecell[{{p{\linewidth}}}]{
                \texttt{\tiny{[P1$\langle$GM]}}
                \texttt{Penalità: (1,6) non è vuoto, ma contiene '─'.} \\
\\ 
\texttt{Avete accumulato collettivamente 15 di 16 penalità.} \\
\texttt{Per favore riprova!} \\
            }
        }
    }
    & \\ \\

    \theutterance \stepcounter{utterance}  
    & \multicolumn{4}{p{0.6\linewidth}}{
        \cellcolor[rgb]{0.9,0.9,0.9}{
            \makecell[{{p{\linewidth}}}]{
                \texttt{\tiny{[P1$\rangle$GM]}}
                \texttt{MUOVERE: M, (7,6)} \\
            }
        }
    }
    & & \\ \\

    \theutterance \stepcounter{utterance}  
    & & & \multicolumn{2}{p{0.3\linewidth}}{
        \cellcolor[rgb]{0.9,0.9,0.9}{
            \makecell[{{p{\linewidth}}}]{
                \texttt{\tiny{[GM$|$GM]}}
                \texttt{MUOVERE: M, (7,6)} \\
            }
        }
    }
    & & \\ \\

    \theutterance \stepcounter{utterance}  
    & & & \multicolumn{2}{p{0.3\linewidth}}{
        \cellcolor[rgb]{0.9,0.9,0.9}{
            \makecell[{{p{\linewidth}}}]{
                \texttt{\tiny{[GM$|$GM]}}
                \texttt{M moved to (7, 6)} \\
            }
        }
    }
    & & \\ \\

    \theutterance \stepcounter{utterance}  
    & & \multicolumn{4}{p{0.6\linewidth}}{
        \cellcolor[rgb]{0.9,0.9,0.9}{
            \makecell[{{p{\linewidth}}}]{
                \texttt{\tiny{[GM$\rangle$P2]}}
                \texttt{Il tuo messaggio è stato trasmesso all'altro giocatore.} \\
\\ 
\\ 
\texttt{Stai giocando il round 9 di un massimo di 28.} \\
\texttt{Avete accumulato collettivamente 15 di 16 penalità.} \\
\texttt{L'altro giocatore ha spostato un oggetto sulla sua griglia.} \\
\\ 
\\ 
\texttt{Qual è il tuo prossimo comando?} \\
            }
        }
    }
    & \\ \\

    \theutterance \stepcounter{utterance}  
    & & & \multicolumn{4}{p{0.6\linewidth}}{
        \cellcolor[rgb]{0.9,0.9,0.9}{
            \makecell[{{p{\linewidth}}}]{
                \texttt{\tiny{[GM$\langle$P2]}}
                \texttt{Siamo a 15 penalità su 16, devo stare molto attento. Aspetto la risposta dell'altro giocatore prima di muovere M, ma intanto posso spostare **S** in (7,3) come proposto prima. Controllo: riga 3, colonna 7 è `─`, non vuoto. Quindi non posso. Devo invece spostare **P** in (4,5) come precedentemente proposto. Controllo: riga 5, colonna 4 è `│`, non vuoto. Allora propongo nuove alternative.} \\
\\ 
\texttt{DIRE: S a (7,3) e P a (4,5) da me non sono vuoti. Propongo S a (7,2) e P a (4,6). Per M attendo tua conferma su (2,6).} \\
            }
        }
    }
     \\ \\

    \theutterance \stepcounter{utterance}  
    & & & \multicolumn{2}{p{0.3\linewidth}}{
        \cellcolor[rgb]{0.9,0.9,0.9}{
            \makecell[{{p{\linewidth}}}]{
                \texttt{\tiny{[GM$|$GM]}}
                \texttt{Siamo a 15 penalità su 16, devo stare molto attento. Aspetto la risposta dell'altro giocatore prima di muovere M, ma intanto posso spostare **S** in (7,3) come proposto prima. Controllo: riga 3, colonna 7 è `─`, non vuoto. Quindi non posso. Devo invece spostare **P** in (4,5) come precedentemente proposto. Controllo: riga 5, colonna 4 è `│`, non vuoto. Allora propongo nuove alternative.} \\
\\ 
\texttt{DIRE: S a (7,3) e P a (4,5) da me non sono vuoti. Propongo S a (7,2) e P a (4,6). Per M attendo tua conferma su (2,6).} \\
            }
        }
    }
    & & \\ \\

    \theutterance \stepcounter{utterance}  
    & & & \multicolumn{2}{p{0.3\linewidth}}{
        \cellcolor[rgb]{0.9,0.9,0.9}{
            \makecell[{{p{\linewidth}}}]{
                \texttt{\tiny{[GM$|$GM]}}
                \texttt{Invalid format: head is not empty: 'Siamo a 15 penalità su 16, devo stare molto attento. Aspetto la risposta dell'altro giocatore prima di muovere M, ma intanto posso spostare **S** in (7,3) come proposto prima. Controllo: riga 3, colonna 7 è ─, non vuoto. Quindi non posso. Devo invece spostare **P** in (4,5) come precedentemente proposto. Controllo: riga 5, colonna 4 è │, non vuoto. Allora propongo nuove alternative.} \\
\\ 
\texttt{'} \\
            }
        }
    }
    & & \\ \\

    \theutterance \stepcounter{utterance}  
    & & \multicolumn{4}{p{0.6\linewidth}}{
        \cellcolor[rgb]{0.9,0.9,0.9}{
            \makecell[{{p{\linewidth}}}]{
                \texttt{\tiny{[GM$\rangle$P2]}}
                \texttt{Penalità: Il tuo messaggio non deve contenere nulla prima del comando!} \\
\\ 
\texttt{Assicurati che la tua risposta contenga solo `DIRE: <MESSAGGIO>` o `MUOVERE: <OGGETTO>, (<X>, <Y>)`, e nient'altro!} \\
\\ 
\texttt{Avete accumulato collettivamente 16 di 16 penalità.} \\
\texttt{Per favore riprova!} \\
            }
        }
    }
    & \\ \\

    \theutterance \stepcounter{utterance}  
    & & & \multicolumn{4}{p{0.6\linewidth}}{
        \cellcolor[rgb]{0.9,0.9,0.9}{
            \makecell[{{p{\linewidth}}}]{
                \texttt{\tiny{[GM$\langle$P2]}}
                \texttt{DIRE: S a (7,3) e P a (4,5) da me non sono vuoti. Propongo S a (7,2) e P a (4,6). Per M attendo tua conferma su (2,6).} \\
            }
        }
    }
     \\ \\

    \theutterance \stepcounter{utterance}  
    & & & \multicolumn{2}{p{0.3\linewidth}}{
        \cellcolor[rgb]{0.9,0.9,0.9}{
            \makecell[{{p{\linewidth}}}]{
                \texttt{\tiny{[GM$|$GM]}}
                \texttt{DIRE: S a (7,3) e P a (4,5) da me non sono vuoti. Propongo S a (7,2) e P a (4,6). Per M attendo tua conferma su (2,6).} \\
            }
        }
    }
    & & \\ \\

    \theutterance \stepcounter{utterance}  
    & & \multicolumn{4}{p{0.6\linewidth}}{
        \cellcolor[rgb]{0.9,0.9,0.9}{
            \makecell[{{p{\linewidth}}}]{
                \texttt{\tiny{[P1$\langle$GM]}}
                \texttt{Spostato 'M' in (7,6) con successo. La tua griglia aggiornata appare così:} \\
\texttt{```} \\
\texttt{ 1234567} \\
\texttt{╔═══════╗} \\
\texttt{║◌◌◌◌◌◌T║ 1} \\
\texttt{║P◌O◌◌◌◌║ 2} \\
\texttt{║◌◌◌◌◌┌─╢ 3} \\
\texttt{║◌D◌┌─┤S║ 4} \\
\texttt{║◌◌A│◌├─╢ 5} \\
\texttt{╟───┼─┘M║ 6} \\
\texttt{║◌◌◌│◌◌◌║ 7} \\
\texttt{╚═══╧═══╝} \\
\texttt{```} \\
\texttt{Stai giocando il round 9 di un massimo di 28.} \\
\texttt{Avete accumulato collettivamente 16 di 16 penalità.} \\
\texttt{L'altro giocatore ha inviato questo messaggio:} \\
\\ 
\\ 
\texttt{"S a (7,3) e P a (4,5) da me non sono vuoti. Propongo S a (7,2) e P a (4,6). Per M attendo tua conferma su (2,6)."} \\
\\ 
\\ 
\texttt{Qual è il tuo prossimo comando?} \\
            }
        }
    }
    & \\ \\

    \theutterance \stepcounter{utterance}  
    & \multicolumn{4}{p{0.6\linewidth}}{
        \cellcolor[rgb]{0.9,0.9,0.9}{
            \makecell[{{p{\linewidth}}}]{
                \texttt{\tiny{[P1$\rangle$GM]}}
                \texttt{DIRE: Attenzione! Abbiamo raggiunto 16 penalità. Dobbiamo evitare altri errori. M l'ho spostato a (7,6) perché (2,6) e (1,6) contengono linee. S a (7,2) per me va bene. P a (4,6) da me ha una linea però, controlla. Propongo P a (5,5) che è vuoto.} \\
            }
        }
    }
    & & \\ \\

    \theutterance \stepcounter{utterance}  
    & & & \multicolumn{2}{p{0.3\linewidth}}{
        \cellcolor[rgb]{0.9,0.9,0.9}{
            \makecell[{{p{\linewidth}}}]{
                \texttt{\tiny{[GM$|$GM]}}
                \texttt{DIRE: Attenzione! Abbiamo raggiunto 16 penalità. Dobbiamo evitare altri errori. M l'ho spostato a (7,6) perché (2,6) e (1,6) contengono linee. S a (7,2) per me va bene. P a (4,6) da me ha una linea però, controlla. Propongo P a (5,5) che è vuoto.} \\
            }
        }
    }
    & & \\ \\

    \theutterance \stepcounter{utterance}  
    & & \multicolumn{4}{p{0.6\linewidth}}{
        \cellcolor[rgb]{0.9,0.9,0.9}{
            \makecell[{{p{\linewidth}}}]{
                \texttt{\tiny{[GM$\rangle$P2]}}
                \texttt{Il tuo messaggio è stato trasmesso all'altro giocatore.} \\
\\ 
\\ 
\texttt{Stai giocando il round 10 di un massimo di 28.} \\
\texttt{Avete accumulato collettivamente 16 di 16 penalità.} \\
\texttt{L'altro giocatore ha inviato questo messaggio:} \\
\\ 
\\ 
\texttt{"Attenzione! Abbiamo raggiunto 16 penalità. Dobbiamo evitare altri errori. M l'ho spostato a (7,6) perché (2,6) e (1,6) contengono linee. S a (7,2) per me va bene. P a (4,6) da me ha una linea però, controlla. Propongo P a (5,5) che è vuoto."} \\
\\ 
\\ 
\texttt{Qual è il tuo prossimo comando?} \\
            }
        }
    }
    & \\ \\

    \theutterance \stepcounter{utterance}  
    & & & \multicolumn{4}{p{0.6\linewidth}}{
        \cellcolor[rgb]{0.9,0.9,0.9}{
            \makecell[{{p{\linewidth}}}]{
                \texttt{\tiny{[GM$\langle$P2]}}
                \texttt{Controllo la mia griglia:} \\
\texttt{{-} M a (7,6): riga 6, colonna 7 è `◌` – vuoto, OK.} \\
\texttt{{-} S a (7,2): riga 2, colonna 7 è `◌` – vuoto, OK.} \\
\texttt{{-} P a (5,5): riga 5, colonna 5 è `◌` – vuoto, OK.} \\
\\ 
\texttt{Posso iniziare a spostare. Sposto S in (7,2).} \\
\\ 
\texttt{MUOVERE: S, (7, 2)} \\
            }
        }
    }
     \\ \\

    \theutterance \stepcounter{utterance}  
    & & & \multicolumn{2}{p{0.3\linewidth}}{
        \cellcolor[rgb]{0.9,0.9,0.9}{
            \makecell[{{p{\linewidth}}}]{
                \texttt{\tiny{[GM$|$GM]}}
                \texttt{Controllo la mia griglia:} \\
\texttt{{-} M a (7,6): riga 6, colonna 7 è `◌` – vuoto, OK.} \\
\texttt{{-} S a (7,2): riga 2, colonna 7 è `◌` – vuoto, OK.} \\
\texttt{{-} P a (5,5): riga 5, colonna 5 è `◌` – vuoto, OK.} \\
\\ 
\texttt{Posso iniziare a spostare. Sposto S in (7,2).} \\
\\ 
\texttt{MUOVERE: S, (7, 2)} \\
            }
        }
    }
    & & \\ \\

    \theutterance \stepcounter{utterance}  
    & & & \multicolumn{2}{p{0.3\linewidth}}{
        \cellcolor[rgb]{0.9,0.9,0.9}{
            \makecell[{{p{\linewidth}}}]{
                \texttt{\tiny{[GM$|$GM]}}
                \texttt{Invalid format: head is not empty: 'Controllo la mia griglia:} \\
\texttt{{-} M a (7,6): riga 6, colonna 7 è ◌ – vuoto, OK.} \\
\texttt{{-} S a (7,2): riga 2, colonna 7 è ◌ – vuoto, OK.} \\
\texttt{{-} P a (5,5): riga 5, colonna 5 è ◌ – vuoto, OK.} \\
\\ 
\texttt{Posso iniziare a spostare. Sposto S in (7,2).} \\
\\ 
\texttt{'} \\
            }
        }
    }
    & & \\ \\

    \theutterance \stepcounter{utterance}  
    & & & \multicolumn{2}{p{0.3\linewidth}}{
        \cellcolor[rgb]{0.9,0.9,0.9}{
            \makecell[{{p{\linewidth}}}]{
                \texttt{\tiny{[GM$|$GM]}}
                \texttt{Maximum number of penalties exceeded} \\
            }
        }
    }
    & & \\ \\

    \theutterance \stepcounter{utterance}  
    & & & \multicolumn{2}{p{0.3\linewidth}}{
        \cellcolor[rgb]{0.9,0.9,0.9}{
            \makecell[{{p{\linewidth}}}]{
                \texttt{\tiny{[GM$|$GM]}}
                \texttt{Initial states:} \\
\texttt{```} \\
\texttt{Player 1:} \\
\texttt{ 1234567} \\
\texttt{╔═══════╗} \\
\texttt{║◌◌◌◌◌T◌║ 1} \\
\texttt{║P◌◌◌◌◌◌║ 2} \\
\texttt{║◌◌AD◌┌─╢ 3} \\
\texttt{║◌◌◌┌─┤S║ 4} \\
\texttt{║O◌◌│◌├─╢ 5} \\
\texttt{╟───┼─┘◌║ 6} \\
\texttt{║◌M◌│◌◌◌║ 7} \\
\texttt{╚═══╧═══╝} \\
\\ 
\\ 
\texttt{Player 2:} \\
\texttt{ 1234567} \\
\texttt{╔═══════╗} \\
\texttt{║◌◌◌◌◌S◌║ 1} \\
\texttt{║◌◌◌◌O◌◌║ 2} \\
\texttt{║◌◌◌M◌┌─╢ 3} \\
\texttt{║◌◌A┌─┤◌║ 4} \\
\texttt{║◌◌◌│◌├─╢ 5} \\
\texttt{╟───┼─┘◌║ 6} \\
\texttt{║D◌T│P◌◌║ 7} \\
\texttt{╚═══╧═══╝} \\
\\ 
\texttt{```} \\
            }
        }
    }
    & & \\ \\

    \theutterance \stepcounter{utterance}  
    & & & \multicolumn{2}{p{0.3\linewidth}}{
        \cellcolor[rgb]{0.9,0.9,0.9}{
            \makecell[{{p{\linewidth}}}]{
                \texttt{\tiny{[GM$|$GM]}}
                \texttt{End states:} \\
\texttt{```} \\
\texttt{Player 1:} \\
\texttt{ 1234567} \\
\texttt{╔═══════╗} \\
\texttt{║◌◌◌◌◌◌T║ 1} \\
\texttt{║P◌O◌◌◌◌║ 2} \\
\texttt{║◌◌◌◌◌┌─╢ 3} \\
\texttt{║◌D◌┌─┤S║ 4} \\
\texttt{║◌◌A│◌├─╢ 5} \\
\texttt{╟───┼─┘M║ 6} \\
\texttt{║◌◌◌│◌◌◌║ 7} \\
\texttt{╚═══╧═══╝} \\
\\ 
\\ 
\texttt{Player 2:} \\
\texttt{ 1234567} \\
\texttt{╔═══════╗} \\
\texttt{║◌◌◌◌◌S◌║ 1} \\
\texttt{║◌◌O◌◌◌◌║ 2} \\
\texttt{║◌◌◌M◌┌─╢ 3} \\
\texttt{║◌D◌┌─┤T║ 4} \\
\texttt{║◌◌A│◌├─╢ 5} \\
\texttt{╟───┼─┘◌║ 6} \\
\texttt{║◌◌◌│P◌◌║ 7} \\
\texttt{╚═══╧═══╝} \\
\\ 
\texttt{```} \\
            }
        }
    }
    & & \\ \\

    \theutterance \stepcounter{utterance}  
    & & & \multicolumn{2}{p{0.3\linewidth}}{
        \cellcolor[rgb]{0.9,0.9,0.9}{
            \makecell[{{p{\linewidth}}}]{
                \texttt{\tiny{[GM$|$GM]}}
                \texttt{* success: False} \\
\texttt{* lose: False} \\
\texttt{* aborted: True} \\
\texttt{{-}{-}{-}{-}{-}{-}{-}} \\
\texttt{* Shifts: 6.00} \\
\texttt{* Max Shifts: 12.00} \\
\texttt{* Min Shifts: 6.00} \\
\texttt{* End Distance Sum: 16.81} \\
\texttt{* Init Distance Sum: 31.75} \\
\texttt{* Expected Distance Sum: 29.33} \\
\texttt{* Penalties: 17.00} \\
\texttt{* Max Penalties: 16.00} \\
\texttt{* Rounds: 9.00} \\
\texttt{* Max Rounds: 28.00} \\
\texttt{* Object Count: 7.00} \\
            }
        }
    }
    & & \\ \\

\end{supertabular}
}

\end{document}
